\documentclass[11pt]{article}
\usepackage[T1]{fontenc}
\usepackage[utf8]{inputenc}
\usepackage{lmodern}
\usepackage[margin=1in]{geometry}
\usepackage{amsmath,amssymb,amsthm,mathtools}
\usepackage{aliascnt}
\usepackage{microtype}
\usepackage{booktabs,array,longtable,tabularx}
\usepackage{xcolor}
\usepackage{enumitem}
\usepackage{fancyhdr}
\usepackage{needspace}
\usepackage{hyperref}
\usepackage[nameinlink,noabbrev]{cleveref}
\hypersetup{colorlinks=true,linkcolor=blue!45!black,citecolor=blue!45!black,
 urlcolor=blue!45!black,
 pdftitle={Cofinality, Factorization and Scalar Extension for Entire Hahn Functions at Arbitrary Rank},
 pdfauthor={Merged research report prepared for Vladimir Reshetnikov}}
\setlength{\headheight}{14pt}
\pagestyle{fancy}\fancyhf{}
\fancyhead[L]{\small Entire Hahn functions at arbitrary rank}
\fancyhead[R]{\small Cofinality, B\'ezout failure, scalar extension}
\fancyfoot[C]{\thepage}
\setlength{\parskip}{3pt}
\setcounter{tocdepth}{1}
\widowpenalty=10000\clubpenalty=10000
\setlength{\emergencystretch}{2em}
\allowdisplaybreaks[2]
\newtheorem{theorem}{Theorem}[section]
\theoremstyle{plain}
\newaliascnt{lemma}{theorem}
\newtheorem{lemma}[lemma]{Lemma}
\aliascntresetthe{lemma}
\crefname{lemma}{lemma}{lemmas}
\Crefname{lemma}{Lemma}{Lemmas}
\theoremstyle{plain}
\newaliascnt{proposition}{theorem}
\newtheorem{proposition}[proposition]{Proposition}
\aliascntresetthe{proposition}
\crefname{proposition}{proposition}{propositions}
\Crefname{proposition}{Proposition}{Propositions}
\theoremstyle{plain}
\newaliascnt{corollary}{theorem}
\newtheorem{corollary}[corollary]{Corollary}
\aliascntresetthe{corollary}
\crefname{corollary}{corollary}{corollaries}
\Crefname{corollary}{Corollary}{Corollaries}
\theoremstyle{definition}
\newaliascnt{definition}{theorem}
\newtheorem{definition}[definition]{Definition}
\aliascntresetthe{definition}
\crefname{definition}{definition}{definitions}
\Crefname{definition}{Definition}{Definitions}
\theoremstyle{definition}
\newaliascnt{convention}{theorem}
\newtheorem{convention}[convention]{Convention}
\aliascntresetthe{convention}
\crefname{convention}{convention}{conventions}
\Crefname{convention}{Convention}{Conventions}
\theoremstyle{definition}
\newaliascnt{example}{theorem}
\newtheorem{example}[example]{Example}
\aliascntresetthe{example}
\crefname{example}{example}{examples}
\Crefname{example}{Example}{Examples}
\theoremstyle{definition}
\newaliascnt{question}{theorem}
\newtheorem{question}[question]{Question}
\aliascntresetthe{question}
\crefname{question}{question}{questions}
\Crefname{question}{Question}{Questions}
\theoremstyle{remark}
\newaliascnt{remark}{theorem}
\newtheorem{remark}[remark]{Remark}
\aliascntresetthe{remark}
\crefname{remark}{remark}{remarks}
\Crefname{remark}{Remark}{Remarks}
\newcommand{\C}{\mathbb C}\newcommand{\R}{\mathbb R}
\newcommand{\Q}{\mathbb Q}\newcommand{\Z}{\mathbb Z}\newcommand{\N}{\mathbb N}
\newcommand{\No}{\mathbf{No}}
\newcommand{\K}{K_\Gamma}\newcommand{\KD}{K_\Delta}
\newcommand{\E}{\mathcal E_\Gamma}
\newcommand{\A}{\mathcal A_\Gamma}\newcommand{\OO}{\mathcal O_\Gamma}
\newcommand{\TT}{\mathcal T_\Gamma}
\newcommand{\mm}{\mathfrak m_\Gamma}
\newcommand{\Dom}{\mathcal D_{\Gamma,\Delta}}
\DeclareMathOperator{\supp}{supp}\DeclareMathOperator{\lc}{lc}
\DeclareMathOperator{\ord}{ord}\DeclareMathOperator{\cf}{cf}
\DeclareMathOperator{\Div}{Div}\DeclareMathOperator{\inpol}{in}
\DeclareMathOperator{\conv}{conv}
\newcommand{\ggA}{\mathrel{\gg_{\!A}}}
\newcommand{\abs}[1]{\lvert#1\rvert}
\newcommand{\norm}[1]{\lVert#1\rVert}
\newcommand{\res}{\operatorname{res}}
\newcolumntype{Y}{>{\raggedright\arraybackslash}X}

\title{\textbf{Cofinality, Factorization and Scalar Extension\\
for Entire Hahn Functions at Arbitrary Rank}\\[5pt]
\large A classification of the entire ring, its ideal theory,\\
and its behaviour under enlargement of the value group}
\author{Merged research report prepared for Vladimir Reshetnikov}
\date{21 September 2026}

\begin{document}
\maketitle
\begin{abstract}
Let $\Gamma$ be a nonzero set-sized divisible ordered abelian group and
$\K=\C((t^\Gamma))$ the full Hahn field over the trivially valued complex
numbers. We study ordinary power series $\sum a_nZ^n$ whose evaluation
families are strongly Hahn-summable at every point of this one fixed
field. An exact coefficient criterion holds: this happens precisely when
$v(a_n)/n$ tends cofinally to $+\infty$ in $\Gamma$. Consequently
nonpolynomial entire series exist exactly when $\cf(\Gamma)=\aleph_0$.
Support-controlled preparation, proved by well-founded recursion over a
well-ordered monoid rather than by any contraction mapping, yields exact
root counts on every valuation ball, shell and residue direction, a
genus-zero canonical product for every ball-finite divisor with no
exponential correction factors, and complete factorization by zeros at
arbitrary rank.

Two sharp classifications follow. The first is ideal-theoretic. The entire
ring $\E$ is always a GCD domain with units $\K^\times$; it is the
polynomial ring, hence a PID, when $\cf(\Gamma)>\aleph_0$; it is a
non-Noetherian B\'ezout domain with unrestricted Hermite interpolation when
$\Gamma$ has an order unit; and otherwise it is a non-Noetherian GCD domain
that is \emph{not} B\'ezout, witnessed by two explicit entire functions with
disjoint simple zero sets generating a proper ideal. A separate order-unit
criterion decides invertibility in the valuation-restricted algebra
$\TT$, and shows that the positive case of the trichotomy is sharp rather
than merely sufficient. The second classification is analytic. For every
ordered-group extension $\Gamma\subseteq\Delta$, all nonpolynomial entire
series over $\K$ share one exact strong evaluation domain in $\KD$: the
valuation ring of $v_\Delta$ coarsened by the convex hull of $\Gamma$, whose
residue field is $\C((t^H))$ rather than $\C$. Entireness survives exactly
for cofinal extensions, no new zeros appear anywhere in the surviving
domain, and on a noncofinal extension no alternative entire series agrees
with the old one even on all ordinary complex constants. The scale boundary
is intrinsic: the value group, the coefficient slopes and the negated zero
valuations generate the same convex subgroup.

This report is the union of two independently written manuscripts that
proved the same engine. The arbitrary-rank classifications, the explicit
non-B\'ezout construction and the universal extension domain are proposed
original contributions, not a claimed solution of any named published
conjecture. Classical Hahn-field inputs and rank-one antecedents are
identified and credited. Complete proofs are supplied; they have not been
independently refereed or verified in a proof assistant.
\end{abstract}
\newpage
\tableofcontents
\newpage

\section{The question, the answer, and its scope}
\label{ent:sec:intro}

\subsection{Two questions that turn out to have one engine}
Over an algebraically closed field, finite polynomials with disjoint zero
sets generate the unit ideal. In familiar one-variable analytic theories,
interpolation extends this fact to entire functions, and it is tempting to
infer the same conclusion as soon as an infinite product theorem is
available. The inference is false in the higher-rank setting studied here.
We obtain a ring in which every admissible zero divisor is realized, every
function factors completely according to its zeros, and greatest common
divisors exist, yet two functions without a common zero need not satisfy a
B\'ezout identity.

The second question is what the word \emph{entire} names. A formal series
with coefficients in one Hahn field may be evaluable at every point of that
field and lose even strong summability at a single point of a larger one.
Both questions are decided by the same object: the ordered hierarchy of
scales above the coefficient slopes of the series. The first is decided by
whether those scales have a greatest Archimedean class; the second by
whether they remain cofinal after the enlargement.

The obstruction behind the first question is genuinely infinite. At
increasingly large arguments, any one entire function is restricted by a
common lower bound on the valuations of its coefficients. We place a second
sequence of zeros so close to the first that the reciprocal values required
by a B\'ezout identity escape those bounds. Every finite interpolation
problem remains solvable by a polynomial. It is the simultaneous infinite
problem that fails.

\subsection{The analytic category}
Throughout, $\Gamma$ is a \emph{set-sized divisible ordered abelian group},
written additively, and
\[
 \K=\C((t^\Gamma)),\qquad
 v\left(\sum_\eta c_\eta t^\eta\right)=\min\{\eta:c_\eta\ne0\}.
\]
A Hahn series has well-ordered support. We put $v(0)=+\infty$. This is the
\emph{full} Hahn field: all well-ordered supports are admitted, not a
support-restricted subfield. The complex coefficient field carries the
\emph{trivial} valuation, so every nonzero complex coefficient has valuation
zero and an infinite ordinary complex sum at one exponent is not an allowed
Hahn coefficient operation. The valuation ring and its maximal ideal are
\[
 \OO=\{x\in\K:v(x)\ge0\},\qquad
 \mm=\{x\in\K:v(x)>0\},
\]
with residue field exactly $\C$.

\begin{definition}\label{ent:def:entire}
The ring of \emph{strongly entire Hahn functions} is
\[
 \E=\left\{F=\sum_{n\ge0}a_nZ^n\in\K[[Z]]:
       (a_nz^n)_{n\ge0}\text{ is strongly Hahn-summable for every }z\in\K
       \right\}.
\]
Here strong summability means that the union of the term supports is well
ordered and each exponent belongs to only finitely many term supports. That
this set is a ring is proved below.
\end{definition}

We use $\N=\{0,1,2,\ldots\}$. The variable $Z$ is an ordinary single
indeterminate; its exponents are nonnegative \emph{integers}. The Hahn
exponents occur only in the coefficients, and the two infinities must not be
confused: a finite index $n$ in a power series is not a Hahn exponent, and
integer multiples of a positive Hahn exponent need not be cofinal. Nothing
is asserted about arbitrary generalized series in $Z$ with ordinal support
order types, about arbitrary local functions for the fine surreal topology,
or about the Berarducci--Mantova derivation. All derivatives in this report
hold $\K$ constant.

\begin{convention}[Ring discipline]\label{ent:conv:ring}
Every numbered statement below is asserted over one named coefficient ring,
and no statement is transferred to another. The rings used here are
\[
 \K[Z]\ \subseteq\ \E\ \subseteq\ \TT\ \subseteq\ \A\ \subseteq\ \K[[Z]],
\]
defined in \cref{ent:sec:support,ent:sec:units}, all of which have
\emph{Hahn series over $\C$} as coefficients. In particular no theorem here
is a theorem about $\mathcal O(U)((t^\Gamma))$ or $\C\{z\}((t^\Gamma))$,
whose coefficients are ordinary holomorphic functions on a common complex
domain or convergent germs. That distinction is not a formality: the
companion report \emph{Analytic Geometry over Surcomplex Hahn Fields}
adopts the same discipline by name for its own four coefficient rings, and
records there an explicit element of the radius-free algebra lying in no
common-domain algebra. We invoke that convention here so that no reader
transports a statement of this report into those categories. See
\cref{ent:sec:collection}.
\end{convention}

\subsection{The two classifications}
A subset of $\Gamma$ is cofinal if every element of $\Gamma$ is at most some
member of that subset; $\cf(\Gamma)$ is the least cardinality of a cofinal
subset, and it is at least $\aleph_0$ for $\Gamma\ne0$, since an ordered
group has no greatest element.

\begin{definition}[Order unit]\label{ent:def:order-unit}
An element $\delta>0$ of $\Gamma$ is an \emph{order unit} if for every
$\gamma\in\Gamma$ some positive integer $m$ satisfies $m\delta>\gamma$; that
is, $\{m\delta:m\ge1\}$ is cofinal in $\Gamma$.
\end{definition}

\begin{remark}[One name for one invariant]\label{ent:rem:one-name}
This report uses exactly two order-theoretic invariants of $\Gamma$, and
they must not be conflated.
\begin{itemize}[leftmargin=1.4em]
\item \emph{Cofinality} always means $\cf(\Gamma)$, the least cardinality of
a cofinal subset of the underlying ordered \emph{set} of the value group. It
is never the cofinality of an ordinal index, of a diagram, or of the image of
some displacement or spectral map; other reports in this collection use the
word for such other invariants, and nothing here transfers to them.
\item \emph{Order unit} is the condition of
\cref{ent:def:order-unit}. The two source manuscripts merged here named it
twice: one called $\delta$ an order unit, the other called it a
\emph{cofinal-cyclic element}, meaning exactly that the cyclic subgroup it
generates is cofinal. These are the same condition, stated twice; we keep
the single name \emph{order unit} and record the synonym here once.
It is equivalent to each of the following: the positive Archimedean classes
of $\Gamma$ have a greatest member; $m\delta\to+\infty$; the monomial
$t^\delta$ is topologically nilpotent in the intrinsic valuation topology of
$\K$.
\end{itemize}
An order unit forces $\cf(\Gamma)=\aleph_0$, but not conversely: the group
$\Gamma_\infty$ of \cref{ent:sec:examples} has countable cofinality and no
order unit. That gap is exactly the third case of the trichotomy below, and
also the difference between the two existence questions of
\cref{ent:thm:main,ent:thm:units}.
\end{remark}

\Needspace{18\baselineskip}
\begin{theorem}[Cofinality--rank trichotomy]\label{ent:thm:main}
Let $\Gamma\ne0$ be a divisible ordered abelian group. Then $\E$ is a GCD
domain, its units are exactly $\K^\times$, and exactly one of the following
cases occurs.
\begin{enumerate}[label=\textnormal{(\roman*)},leftmargin=2em]
\item If $\cf(\Gamma)>\aleph_0$, then $\E=\K[Z]$, a principal ideal domain.
\item If $\Gamma$ has an order unit, then $\E$ contains nonpolynomial
functions and is a non-Noetherian B\'ezout domain. Arbitrary finite-order
jets on every radially finite set can be interpolated by an entire function.
\item If $\cf(\Gamma)=\aleph_0$ and $\Gamma$ has no order unit, then $\E$
contains nonpolynomial functions and is a non-Noetherian GCD domain that is
not B\'ezout. More strongly, there are $f,g\in\E$ with disjoint simple zero
sets and $(f,g)\ne\E$.
\end{enumerate}
For $\Gamma=0$, the strong-summation convention gives $\mathcal E_0=\C[Z]$.
\end{theorem}

A GCD domain is an integral domain in which every pair has a greatest common
divisor, defined up to a unit. A B\'ezout domain is one in which every
finitely generated ideal is principal. Case (iii) therefore separates two
properties that coincide for finite polynomial algebra. This is not a
distinction about algebraic closedness: every $\K$ here is algebraically
closed.

\Needspace{14\baselineskip}
\begin{theorem}[Universal scalar-extension package]\label{ent:thm:main-extension}
Let $\Gamma\subseteq\Delta$ be nonzero divisible ordered abelian groups and
let $H=\conv_\Delta(\Gamma)$ be the convex hull of $\Gamma$ in $\Delta$.
Write $\Dom$ for the valuation ring of $v_\Delta$ coarsened by $H$. Then for
every \emph{nonpolynomial} $f\in\E$:
\begin{enumerate}[label=\textnormal{(\alph*)},leftmargin=2em]
\item the family $(a_nz^n)_n$ is strongly summable in $\KD$ exactly for
$z\in\Dom$, so the domain does not depend on $f$, and $f(\Dom)\subseteq\Dom$;
\item $\Dom=\KD$, equivalently $f\in\mathcal E_\Delta$, exactly when
$\Gamma$ is cofinal in $\Delta$;
\item every zero of the extended function on $\Dom$ already lies in $\K$
with unchanged multiplicity, even when $f$ has infinitely many zeros;
\item if $\Gamma$ is not cofinal in $\Delta$, there is no
$g\in\mathcal E_\Delta$ agreeing with $f$ on $\C$, let alone on $\K$;
\item $\Gamma$, the slope set $\{v(a_n)/n\}$ and the negated zero valuations
$\{-v(r)\}$ generate the same convex subgroup of any larger ordered group,
so $\Dom$ is recoverable from the coefficients alone or from the zeros alone.
\end{enumerate}
The residue field of $\Dom$ is canonically $\C((t^H))$, not $\C$.
\end{theorem}

\begin{remark}[The two extension statements are one statement]
\label{ent:rem:conflict}
Part (b) is often phrased negatively --- in the noncofinal case the series
is simply \emph{not entire over} $\KD$, and indeed
$\mathcal E_\Delta\cap\K[[Z]]=\K[Z]$ there
(\cref{ent:thm:extension-criterion}) --- while parts (a) and (c) say that
the very same series remains strongly summable on all of $\Dom$ and keeps
every one of its zeros there. Read side by side these look like a
disagreement about whether the function survives an enlargement. They are
not in conflict: they are the coarse and the exact form of one fact. ``Not
entire over $\KD$'' means precisely that the strong evaluation domain is a
\emph{proper} subring of $\KD$, and \cref{ent:thm:extension} computes
exactly which subring it is, namely $\Dom$, on which the function is
unchanged. The negative statement locates the boundary; the positive
statement describes everything inside it. Throughout this report ``entire
over $K$'' is reserved for strong summability at \emph{every} point of $K$,
and any weaker survival is described by naming its domain.
\end{remark}

The proof of \cref{ent:thm:main} is distributed through the article. The
coefficient criterion and cofinality dichotomy are in
\cref{ent:sec:growth}; the order-unit criterion for the restricted algebra,
which is what makes case (ii) sharp rather than merely sufficient, is in
\cref{ent:sec:units}; preparation and root counting are in
\cref{ent:sec:prep,ent:sec:zeros}; factorization and GCDs are in
\cref{ent:sec:products}; the negative and positive B\'ezout results are in
\cref{ent:sec:obstruction,ent:sec:positive}.
\Cref{ent:sec:classification} assembles them.
\Cref{ent:thm:main-extension} is proved in \cref{ent:sec:extension}.

\subsection{Relationship to the repository and the literature}
\label{ent:subsec:repo}
The repository supplied with the question already contains a rank-one report
over $\C((t^\R))$, including preparation, zero counting, zero-free rigidity
and a partial-theta canonical product \cite{RepoRankOne}. Neither source
manuscript claims to originate those points, and neither do we; the present
report is their arbitrary-rank generalization, and the rank-one report
itself proves the obstruction that makes rank one non-generic. The
repository also contains a global-divisor report in a different ring,
$\mathcal O(U)((t^\Gamma))$, with ordinary holomorphic coefficient functions
on a common complex domain \cite{RepoDivisors}; its global support
obstructions belong to a different analytic category and our argument is not
a restatement of its Picard-group or moving-divisor results.
\Cref{ent:sec:collection} states the cross-references in full.

Classical non-Archimedean factorization is an essential antecedent. Lazard's
work is a foundational reference for one-variable zeros \cite{Lazard};
Cherry records the one-variable product factorization and develops GCD and
factorization results under a complete real-valued non-Archimedean absolute
value \cite{Cherry,CherryLectures}, and himself calls those factorization
results well known. Hahn-field algebraic closedness is standard, used here
in the form of Poonen's Corollary~4 \cite{Poonen}; the support calculus is
Neumann's \cite{Neumann}; support-based operator methods are developed
extensively by van der Hoeven \cite{vdH}. The distinction between positive
valuation and topological nilpotence in rank two is established, and we
follow Conrad's discussion of it \cite{Conrad}. We claim none of these facts
as new.

\paragraph{Novelty statement.}
The proposed contributions are the arbitrary-rank coefficient, unit and
workspace criteria, the explicit non-B\'ezout construction, the universal
scalar-extension domain with zero conservation, and their combination into
\cref{ent:thm:main,ent:thm:main-extension}. The support-controlled
preparation and product arguments provide a self-contained route to those
results; their rank-one specializations are not new. The literature checks
performed for the two source manuscripts did not locate the stated
arbitrary-rank classifications or the explicit obstruction in the primary
sources consulted. This is a bounded literature statement, not a
certification of priority. The repository review was targeted at the
catalogue, relevant READMEs and stated scope, not a line-by-line audit of
every manuscript. No named published open problem is claimed solved.
\Cref{ent:sec:scope,ent:app:sources} state the limits in full; every
limitation recorded by either source manuscript is preserved there.

\section{Hahn supports, the surreal interpretation, and the algebra chain}
\label{ent:sec:support}

\subsection{Foundational setting}
The proofs can be carried out in ZFC using only the set-sized field $\K$. To
interpret them inside the surcomplex numbers, take $\Gamma$ to be a divisible
ordered additive subgroup of $\No$ and send
\begin{equation}\label{ent:eq:embedding}
 \sum_\gamma c_\gamma t^\gamma
 \longmapsto \sum_\gamma c_\gamma\omega^{-\gamma}.
\end{equation}
The convention $t^\gamma:=\omega^{-\gamma}$ is Conway's monomial map; it is
explicitly \emph{not} a surreal exponential. The real and imaginary parts of
the image are Conway normal forms, and the standard normal-form arithmetic
makes \eqref{ent:eq:embedding} an embedding of fields: the real coefficient
subfield $\R((t^\Gamma))$ lies in $\No$, and its coefficientwise
complexification lies in $\No[i]$ \cite{Conway,Gonshor}. Conjugation acts on
coefficients only. Our explicit examples use displayed subgroups of $\No$, so
they do not require an unspecified universal embedding theorem. An element
$x\in\OO$ has a unique expression $x=c+\varepsilon$ with $c\in\C$ and
$\varepsilon\in\mm$; the residue map $\res(x)=c$ is defined on $\OO$, not on
all of $\K$.

\subsection{The support facts actually used}
We use the following classical Hahn--Neumann support facts
\cite{Neumann}\cite[Section~3]{Poonen}\cite{vdH}. They are stated explicitly
because every infinite construction below depends on them, and they are
imported, not proved here.

\begin{lemma}[Hahn--Neumann support calculus]\label{ent:lem:neumann}
Let $S,T$ be well-ordered subsets of an ordered abelian group.
\begin{enumerate}[label=\textnormal{(\alph*)},leftmargin=2em]
\item $S\cup T$ and $S+T$ are well ordered. For every $\eta$, there are only
finitely many pairs $(s,t)\in S\times T$ with $s+t=\eta$.
\item If $S\subseteq\Gamma_{>0}$, its additive monoid
$S^*=\{s_1+\cdots+s_r:r\ge0,\ s_i\in S\}$ is well ordered. For each $\eta$,
there are only finitely many finite ordered words in $S$ whose sum is $\eta$.
\item Consequently, if $u$ is a Hahn series over a commutative coefficient
ring and $v(u)>0$, the geometric series $\sum_{r\ge0}u^r$ is strongly
summable and is the inverse of $1-u$.
\end{enumerate}
\end{lemma}

Part (b) is a support theorem, not an assertion that $nv(u)$ tends cofinally
to infinity. The latter can fail. In $\Q e_1\oplus\Q e_2$ with $e_2>ne_1$ for
every $n$, the Hahn identity $\sum_{n\ge0}t^{ne_1}=(1-t^{e_1})^{-1}$ is valid
but its partial sums do not converge in the intrinsic valuation topology: no
tail has valuation greater than $e_2$. \Cref{ent:ex:ranktwo-unit} turns this
observation into a theorem about units.

\begin{lemma}[Cofinal tails]\label{ent:lem:cofinal-tails}
Suppose $u_n\in\K$ and $v(u_n)\to+\infty$ cofinally in $\Gamma$; that is, for
every $B\in\Gamma$, eventually $v(u_n)>B$. Then $(u_n)$ is strongly summable.
Its sum has valuation at least $\min_n v(u_n)$ when at least one term is
nonzero.
\end{lemma}
\begin{proof}
For a fixed $B$, only finitely many nonzero term supports meet
$(-\infty,B]$. A descending sequence in their union, beginning at $B$, would
therefore lie in a finite union of well-ordered supports, which is
impossible. Likewise, any exponent occurs in only finitely many term
supports. The valuation estimate follows from coefficientwise addition. A
nonempty set of term valuations has a minimum because it is a subset of the
well-ordered support union.
\end{proof}

We use the usual ZFC equivalence between well-ordering of a linear order and
the absence of an infinite strictly descending sequence. All supports and
index sets are sets; no proper-class support or proper-class summation is
involved.

\subsection{Three meanings of an infinite expression}
\label{ent:subsec:three-meanings}
It is essential to keep separate
\begin{enumerate}[label=(\arabic*),leftmargin=2em]
\item a strong Hahn sum, which requires only a well-ordered support union
and pointwise finiteness;
\item convergence in the intrinsic valuation topology of the fixed field
$\K$, which requires cofinal progress through every threshold of $\Gamma$;
\item convergence in the topology inherited from all surreal scales.
\end{enumerate}
These need not agree. \Cref{ent:def:entire} starts with (1). The
\emph{all-point} condition will force (2) for evaluation series
(\cref{ent:thm:criterion}), even though (1) alone does not, and on a single
disk the two genuinely differ. We make no assertion of convergence in (3),
and no density of $\C$ in any fine surreal topology is asserted or used.
The repository's foundational reports emphasize this distinction
\cite{RepoCatalogue,RepoRankOne}.

\subsection{The chain of coefficient algebras}
\label{ent:subsec:chain}
Four algebras of series in $Z$ over $\K$ are kept distinct throughout.
Define the \emph{polynomial-coefficient Hahn algebra}
\begin{equation}\label{ent:eq:Halg}
 \A=\C[X]((t^\Gamma)),
\end{equation}
whose elements are $A=\sum_{\eta\in S}P_\eta(X)t^\eta$ with $S$ well ordered
and every $P_\eta$ a finite polynomial, with no common bound on
$\deg P_\eta$; and the \emph{valuation-restricted algebra}
\begin{equation}\label{ent:eq:Talg}
 \TT=\left\{\sum_{n\ge0}a_nZ^n\in\K[[Z]]:
  v(a_n)\to+\infty\text{ cofinally in }\Gamma\right\}.
\end{equation}
Taking the coefficient of each $X^n$ gives an injective ring map
$\A\hookrightarrow\K[[X]]$, and under it $\A$ is exactly the set of series
whose coefficient family $(a_n)_n$ is itself strongly summable: finite degree
at each Hahn exponent is precisely the pointwise-finiteness half of strong
summability. With $X$ and $Z$ identified under that map,
\begin{equation}\label{ent:eq:chain}
 \K[Z]\ \subseteq\ \E\ \subseteq\ \TT\ \subseteq\ \A\ \subseteq\ \K[[Z]].
\end{equation}
The first two inclusions need not be strict; the last two record genuinely
different summation requirements. The two source manuscripts named the ring
\eqref{ent:eq:Halg} differently --- one as the auxiliary algebra $\A$ in the
variable $X$, the other as the strong disk algebra $\mathcal H_\Gamma$ in the
variable $Z$ --- but it is one ring, and we use the single symbol $\A$. That
the inclusions in \eqref{ent:eq:chain} hold, and that $\TT$ and $\A$ are
rings, follows from \cref{ent:lem:neumann} together with
\cref{ent:prop:operations} below.

For $A\ne0$ in $\A$, set $w(A)=\min\{\eta:P_\eta\ne0\}$; the polynomial
$P_{w(A)}$ is its \emph{initial polynomial}. Since $\C[X]$ is a domain,
\begin{equation}\label{ent:eq:gauss-mult}
 w(AB)=w(A)+w(B),
\end{equation}
and initial polynomials multiply. Equivalently, writing
$\inpol_0(A)=\overline{t^{-w(A)}A}\in\C[X]\setminus\{0\}$ for the reduction,
$\inpol_0(AB)=\inpol_0(A)\inpol_0(B)$. This Gauss multiplicativity is the
engine of \cref{ent:sec:units}.

\section{An exact coefficient criterion and the cofinality dichotomy}
\label{ent:sec:growth}

For $n\ge1$, division of $v(a_n)$ by $n$ is legitimate because $\Gamma$ is
divisible and torsion-free. We interpret $(+\infty)/n=+\infty$, so that zero
coefficients impose no condition.

\begin{theorem}[All-point strong summability]\label{ent:thm:criterion}
Let $\Gamma\ne0$ and $F=\sum_{n\ge0}a_nZ^n\in\K[[Z]]$. The following
conditions are equivalent.
\begin{enumerate}[label=\textnormal{(\arabic*)},leftmargin=2em]
\item $F\in\E$.
\item For every $\gamma\in\Gamma$, the family $(a_nt^{n\gamma})_{n\ge0}$ is
strongly summable; equivalently $F(t^\gamma X)\in\A$.
\item For every $h\in\Gamma$, eventually $v(a_n)>nh$.
\item For every $\gamma,B\in\Gamma$, eventually $v(a_n)+n\gamma>B$.
\end{enumerate}
In particular, strong entire evaluation is equivalent to the cofinal
condition $v(a_n)/n\to+\infty$, and it is enough to test the monomial points
$z=t^\gamma$.
\end{theorem}
\begin{proof}
The implication (1)$\Rightarrow$(2) is immediate by testing $z=t^\gamma$; the
reformulation of (2) inside $\A$ holds because the common Hahn support is the
union of the supports of the scaled coefficients, and the coefficient at a
fixed Hahn exponent is a polynomial in $X$ exactly when that exponent occurs
at only finitely many indices $n$. The zero series satisfies all four
conditions, so assume $F\ne0$.

Suppose (2) holds. Applying it at $\gamma=0$ shows that the union of the
supports of the $a_n$ is well ordered. Thus
\[
 m=\min\{v(a_n):a_n\ne0\}\in\Gamma
\]
exists. Suppose (3) fails. There is $h\in\Gamma$ for which infinitely many
positive integers $n$ satisfy $a_n\ne0$ and $v(a_n)\le nh$. Choose $q>0$ with
$q>|h|$ and $q>|m|$, where $|\cdot|$ is the ordered group absolute value.
By (2), the set of leading exponents
\[
 \{v(a_n)-2nq:a_n\ne0\}
\]
is well ordered. Its minimum is attained at some \emph{finite} index $k$;
write that minimum as $\mu$. Then
\[
 \mu=v(a_k)-2kq\ge m-2kq>-(2k+1)q.
\]
Choose a violating index $n>2k+1$. It gives
\[
 v(a_n)-2nq\le nh-2nq<-nq<-(2k+1)q<\mu,
\]
a contradiction. This proves (2)$\Rightarrow$(3). Notice that the argument
used a minimum attained at a finite index; it did not assume that the
multiples of $q$ are cofinal in $\Gamma$. See \cref{ent:app:audit}.

Suppose (3) holds and fix $\gamma,B$. Choose $r>0$ and set
$h=-\gamma+|B|+r$. Eventually
\[
 v(a_n)+n\gamma>n(|B|+r)>B
\]
for $n\ge1$. Thus (4) holds.

Finally, fix $z\ne0$ and put $\gamma=v(z)$. Condition (4) implies
$v(a_nz^n)=v(a_n)+n\gamma\to+\infty$. Apply
\cref{ent:lem:cofinal-tails}. At $z=0$ there is only the constant term.
This proves (4)$\Rightarrow$(1).
\end{proof}

\begin{remark}\label{ent:rem:criterion-strength}
Condition (3) is strictly stronger than $v(a_n)\to+\infty$, that is, than
membership in $\TT$. Over $\Gamma=\Q$ the coefficients $a_n=t^n$ tend to zero
in valuation, but $\sum t^nZ^n$ is not entire: at $Z=t^{-1}$ its terms all
have exponent zero. Conversely $a_n=t^{n^2}$ satisfies (3) over $\Q$ and over
$\R$, but not over a group containing an element larger than every real
number. That last example is exactly the one recorded in the foundations
report of this collection; see \cref{ent:sec:collection}.
\end{remark}

\begin{corollary}[Cofinality dichotomy]\label{ent:cor:cofinality}
For $\Gamma\ne0$, the following are equivalent: $\E\ne\K[Z]$; $\E$ contains a
nonpolynomial series; $\Gamma$ has a countable cofinal subset;
$\cf(\Gamma)=\aleph_0$. If $\cf(\Gamma)>\aleph_0$, every strongly entire
ordinary power series is a polynomial.
\end{corollary}
\begin{proof}
If infinitely many $a_n$ are nonzero, their valuations divided by their
positive indices give, by \cref{ent:thm:criterion}, a countable cofinal
subset of $\Gamma$. Conversely, choose a positive increasing cofinal sequence
$(q_n)_{n\ge1}$. Then
\[
 1+\sum_{n\ge1}t^{nq_n}Z^n
\]
satisfies the coefficient criterion and is nonpolynomial. A nonzero ordered
group has no finite cofinal subset, because it has no greatest element.
\end{proof}

\begin{example}[Two different sources of polynomial rigidity]
\label{ent:ex:two-rigidities}
If $\Gamma=0$, every nonzero coefficient has support $\{0\}$, so strong
summability at $1$ forces a polynomial. If $\Gamma$ has uncountable
cofinality, there are many infinitesimal coefficients, but countably many
coefficient slopes cannot reach all of its scales. The conclusion is again
polynomiality, for a different reason. The same countability argument shows
that in the uncountable-cofinality case even the larger algebra $\TT$
collapses to $\K[Z]$; see \eqref{ent:eq:uncountable-rigidity}.
\end{example}

\begin{proposition}[Entire operations]\label{ent:prop:operations}
$\E$ is a $\K$-subalgebra of $\K[[Z]]$, and so are $\TT$ and $\A$.
Evaluation at each $z\in\K$ is a homomorphism on $\E$. Translation
$F(Z)\mapsto F(Z+b)$, differentiation, and the primitive with zero constant
coefficient preserve $\E$. The translation coefficients are
\begin{equation}\label{ent:eq:taylor}
 [Z^k]F(Z+b)=\sum_{n\ge k}\binom nk a_n b^{n-k}
 =\frac{F^{(k)}(b)}{k!}.
\end{equation}
Entire composition $F\circ G$ is also defined and belongs to $\E$. Moreover,
at each fixed scale $s\in\Gamma$ the partial sums of $F\in\E$ converge to $F$
in the weighted Gauss valuation $m_s$ of \eqref{ent:eq:weighted}.
\end{proposition}
\begin{proof}
The ring assertion for $\E$ follows from \cref{ent:thm:criterion}(2):
addition and multiplication preserve every rescaled copy of $\A$. Explicitly,
at a fixed $s$ write $A_i=v(a_i)+is$ and $B_j=v(b_j)+js$ for two members;
both tend cofinally to infinity and have lower bounds $A,B$. Given $\beta$,
choose $N$ with $A_i>\beta-B$ and $B_i>\beta-A$ for $i>N$; if $n>2N$ then
each pair $i+j=n$ has $i>N$ or $j>N$, so the convolution coefficient has
weighted valuation greater than $\beta$. Taking $s=0$ gives the ring
assertion for $\TT$; the same tail estimate gives the stated weighted
convergence of partial sums. That $\A$ is a ring is
\cref{ent:lem:neumann}(a).

Evaluation can be performed by choosing $\gamma=v(z)$, writing
$z=t^\gamma x$ with $x\in\OO$, and applying \cref{ent:lem:A-units}.

For translation and a fixed testing scale $\gamma$, expand
$F(b+t^\gamma X)$ into its finite binomial blocks. Put
$\beta=\min\{v(b),\gamma\}$, with $v(0)=+\infty$. Every coefficient in the
$n$th block has valuation at least $v(a_n)+n\beta$, which tends cofinally to
infinity, and each block has finitely many $X$-coefficients. The
support-union and local-finiteness proof of \cref{ent:lem:cofinal-tails}
therefore applies to the blocks with polynomial coefficients. This gives an
element of $\A$ at every scale and proves \eqref{ent:eq:taylor}. Regrouping
the same jointly summable family proves the evaluation identity and the
inverse translation by $-b$.

For differentiation, $v(na_n)=v(a_n)$ for positive $n$, since $n$ is a
nonzero complex coefficient of valuation zero. Shifting the index by one
preserves the criterion in \cref{ent:thm:criterion}: for a fixed $h$, choose
$q>0$ with $q>|h|$; the bound $v(a_n)>n(h+q)$ for large $n$ implies the
required bounds with $n-1$ or $n+1$ in place of $n$. No hidden Archimedean
estimate is used. Division of $a_n$ by $n+1$ also preserves valuation, so
termwise integration preserves the criterion as well.

For composition, fix $\gamma$ and let $G_\gamma=G(t^\gamma X)\in\A$. If
$G_\gamma\ne0$, put $\mu=w(G_\gamma)$. Then
$w(a_nG_\gamma^n)=v(a_n)+n\mu\to+\infty$, so the sum is strongly summable in
$\A$. Its coefficients define the same formal composite at every scale, since
rescaling commutes with each coefficient's strong sum. The case $G=0$ is
constant. Evaluation and associativity follow from the same
support-controlled regrouping.
\end{proof}

\section{Restricted series: the sharp order-unit criterion}
\label{ent:sec:units}

The trichotomy of \cref{ent:thm:main} splits its two nonpolynomial cases by
the presence of an order unit. This section proves that the order unit is not
an artefact of the proof technique of \cref{ent:sec:positive}: it is exactly
the condition for a nonconstant element of the restricted algebra $\TT$ of
\eqref{ent:eq:Talg} to be invertible there. That is what makes case (ii) of
the trichotomy sharp rather than merely sufficient, and it is a statement
about $\TT$, not about $\E$; the ring discipline of \cref{ent:conv:ring}
applies.

\begin{theorem}[Restricted-unit criterion]\label{ent:thm:units}
Let $f=\sum a_nZ^n\in\TT$ be nonconstant. Then $f$ is a unit of $\TT$ if and
only if $a_0\ne0$ and
\begin{equation}\label{ent:eq:unit-gap}
 \delta(f)=\min_{n\ge1,\ a_n\ne0}v(a_n/a_0)
\end{equation}
is a positive order unit of $\Gamma$. Nonzero constant series are always
units, and the minimum in \eqref{ent:eq:unit-gap} exists whenever $a_0\ne0$.
\end{theorem}
\begin{proof}
The minimum exists because the coefficient valuations tend cofinally to
infinity: below the valuation of any chosen nonzero coefficient there are
only finitely many nonzero coefficients.

Suppose $fg=1$ with $g\in\TT$. The constant coefficient forces $a_0\ne0$.
Divide $f$ by $a_0$ and multiply $g$ by $a_0$, so that both constant
coefficients are $1$. Then $w(f),w(g)\le0$, while \eqref{ent:eq:gauss-mult}
gives their sum as zero. Hence both are zero. Their reductions in $\C[X]$
multiply to $1$, so both reductions are constant. Every nonconstant
coefficient of $f$ and of $g$ consequently has positive valuation; in
particular $\delta=\delta(f)>0$.

Assume for contradiction that $\delta$ is not an order unit. The convex
subgroup it generates,
\begin{equation}\label{ent:eq:Hdelta}
 H_\delta=\{\gamma\in\Gamma:|\gamma|\le m\delta
                       \text{ for some }m\in\N\},
\end{equation}
is then proper. Coarsen the valuation to the ordered quotient
$\Gamma/H_\delta$. Every coefficient of $f$ and $g$ is integral for this
coarsening. Their coarsened reductions are \emph{polynomials} over the
coarsened residue field: choose $\beta>H_\delta$, meaning $\beta>h$ for every
$h\in H_\delta$; since coefficient valuations tend cofinally to infinity in
$\Gamma$, only finitely many coefficients can have value at most $\beta$, and
all the remaining ones reduce to zero.

The reduced polynomial of $f$ is nonconstant, because a nonconstant
coefficient of valuation $\delta\in H_\delta$ has nonzero coarsened residue.
On the other hand the reduced polynomials of $f$ and $g$ multiply to $1$ in a
polynomial ring over a field, and a nonconstant polynomial cannot be a unit
there. This contradiction proves that $\delta$ is an order unit.

Conversely, suppose $a_0\ne0$ and $\delta(f)$ is a positive order unit.
Write $f=a_0(1+h)$, where $h$ has zero constant coefficient and
$w(h)=\delta$. The formal inverse is
\begin{equation}\label{ent:eq:restricted-inverse}
 f^{-1}=a_0^{-1}\sum_{m\ge0}(-h)^m.
\end{equation}
Its summands have Gauss valuation $m\delta\to+\infty$, so the sum exists
strongly in $\A$. To prove that it belongs to $\TT$, fix a valuation
threshold $\beta$. All sufficiently large $m$ contribute only coefficients of
valuation greater than $\beta$; the remaining finitely many powers of $h$
belong to $\TT$, so their coefficients also eventually exceed $\beta$. Thus
the coefficients of \eqref{ent:eq:restricted-inverse} tend cofinally to
infinity. The finite geometric identity, followed by the justified strong
sum, proves $ff^{-1}=1$.
\end{proof}

The coarsening argument is what makes this a \emph{necessity} theorem rather
than the observation that one particular geometric-series construction fails.
It shows that no cancellation in a hypothetical inverse can evade the
order-unit obstruction.

\begin{corollary}[Unit group of the restricted algebra]
\label{ent:cor:unit-structure}
The algebra $\TT$ has a nonconstant unit if and only if $\Gamma$ has an order
unit; otherwise $\TT^\times=\K^\times$. The rule ``the constant coefficient
strictly dominates all the others'' is a sufficient unit criterion for every
restricted series if and only if $\Gamma$ is Archimedean.
\end{corollary}
\begin{proof}
An order unit $\delta$ makes $1-t^\delta Z$ a nonconstant unit by the
theorem, and the converse follows by applying the theorem to any nonconstant
unit. An ordered group is Archimedean exactly when every positive element has
cofinal positive integer multiples; apply that characterization to
$1-t^\delta Z$ for the last assertion.
\end{proof}

\begin{example}[Zero-free does not mean restricted-invertible]
\label{ent:ex:ranktwo-unit}
Let $\Gamma=\Q\epsilon\oplus\Q\eta$ with $0<\epsilon\ll\eta$, the $\eta$
coordinate being compared first. The polynomial
\[
 f(Z)=1-t^\epsilon Z
\]
has residue $1$ and no zero in $\OO$: its unique zero $t^{-\epsilon}$ has
negative valuation. Yet its only formal inverse,
$\sum_{n\ge0}t^{n\epsilon}Z^n$, is not restricted, since $n\epsilon<\eta$ for
every $n$. It does belong to $\A$ and can be strongly evaluated on the
integral disk. By contrast $1-t^\eta Z$ has a restricted inverse, because
$n\eta$ is cofinal.
\end{example}

\begin{remark}[Why a rank-one contraction proof may not be imported]
\label{ent:rem:no-contraction}
\Cref{ent:ex:ranktwo-unit} is the reason the preparation theorem of
\cref{ent:sec:prep} is proved by well-founded recursion over a well-ordered
monoid and not by a convergent iteration. A rank-one Banach argument
containing the phrase ``the valuation increases by a positive amount at each
step'' is not valid here: the standard iteration may improve the error by
$\epsilon$ repeatedly while remaining forever below $\eta$, and its purported
convergence is then simply unproved. For a valid valuation-convergence proof
the positive amount must be an order unit. The preparation proof below uses
positive \emph{supports} and finite-word control instead, and works with no
order unit at all. The same care applies to \cref{ent:sec:positive}, whose
real-valued coarsening is constructed only after an order unit has been
assumed.
\end{remark}

\begin{remark}[Two existence questions, genuinely different]
\label{ent:rem:two-existence}
Nonpolynomial entire series need only $\cf(\Gamma)=\aleph_0$
(\cref{ent:cor:cofinality}); a nonconstant restricted unit needs an order
unit (\cref{ent:cor:unit-structure}). The group $\Gamma_\infty$ of
\cref{ent:sec:examples} has the first and not the second, so the two
questions separate. In topological language: the valuation topology of $\K$
has a countable neighbourhood base at zero exactly when
$\cf(\Gamma)=\aleph_0$, whereas a nonzero topologically nilpotent element
exists exactly when $\Gamma$ has an order unit, since $v(x^n)=nv(x)$. Neither
of these underlying topological facts is claimed as new; the higher-rank
framework in which they sit is discussed by Conrad \cite{Conrad}.
\end{remark}

\section{Support-controlled preparation}
\label{ent:sec:prep}

\subsection{The algebra that retains support certificates}
Recall $\A=\C[X]((t^\Gamma))$ from \eqref{ent:eq:Halg}. The support facts in
\cref{ent:lem:neumann} justify multiplication in $\A$ and show that the
coefficient map $\A\hookrightarrow\K[[X]]$ respects it.

\begin{lemma}[Units and integral substitution]\label{ent:lem:A-units}
An element of $\A$ is a unit if and only if its initial polynomial is a
nonzero constant. Every $A\in\A$ can be evaluated at any $x\in\OO$, and this
evaluation is a ring homomorphism; if $A\ne0$ the value has valuation at
least $w(A)$. For $x\in\OO$, substitution $X\mapsto X+x$ is an automorphism
of $\A$.
\end{lemma}
\begin{proof}
If the initial polynomial is $c\ne0$ at exponent $\eta$, write
$A=ct^\eta(1+u)$ with $w(u)>0$ and invert $1+u$ using
\cref{ent:lem:neumann}(c). Conversely, if $AB=1$, the product of the two
initial polynomials is $1$, so both are nonzero constants.

For substitution, write $x=c+\varepsilon$ with $c\in\C$ and
$v(\varepsilon)>0$, treating $\varepsilon=0$ separately. For each $\eta$,
expand the finite polynomial $P_\eta(X+c+\varepsilon)$ in powers of
$\varepsilon$. All supports lie in $S+(\supp\varepsilon)^*$, which is well
ordered. At a fixed exponent there are only finitely many decompositions from
$S+(\supp\varepsilon)^*$, finitely many words in $\supp\varepsilon$ for each
such decomposition, and finitely many terms from each $P_\eta$. Thus
substitution is a legitimate element of $\A$. Finite polynomial identities
and strong regrouping prove multiplicativity, and its inverse is substitution
by $-x$. Evaluation follows by setting $X=0$ after this substitution, and the
valuation bound is \cref{ent:lem:cofinal-tails}.
\end{proof}

\subsection{Preparation without an Archimedean contraction}
\begin{theorem}[Support-controlled preparation]\label{ent:thm:prep}
Let $A\in\A$ satisfy $w(A)=0$, and suppose its initial polynomial
$P\in\C[X]$ is monic of degree $d$. There are unique $W,U$ such that
\[
 A=WU,\qquad
 W=P+R,\quad R\in\mm[X],\quad\deg R<d,\qquad
 U=1+Q,\quad w(Q)>0.
\]
Here $W$ is a monic polynomial of degree $d$ and $U$ is a unit of $\A$. If
$S=\supp_t(A-P)\subseteq\Gamma_{>0}$, both corrections are supported in
$S^*\setminus\{0\}$.
\end{theorem}
\begin{proof}
Let $M=S^*$, a well-ordered monoid. Write $A=P+\sum_{\eta>0}A_\eta t^\eta$,
with $A_\eta=0$ off $S$. Recursively for $\eta\in M\setminus\{0\}$, define
polynomials $R_\eta,Q_\eta$ by ordinary division by $P$:
\begin{align}
 H_\eta&=A_\eta-
   \sum_{\substack{\alpha+\beta=\eta\\\alpha,\beta>0}}
       R_\alpha Q_\beta,\label{ent:eq:prep-rec1}\\
 H_\eta&=P Q_\eta+R_\eta,\qquad\deg R_\eta<d.
 \label{ent:eq:prep-rec2}
\end{align}
The sum is finite by \cref{ent:lem:neumann}, and every index on its right is
smaller than $\eta$. Thus this is a well-founded recursion on a set. Set
$R=\sum R_\eta t^\eta$ and $Q=\sum Q_\eta t^\eta$. Because the remainders
have a uniform degree bound, $R$ is a finite polynomial over $\mm$, and
$Q\in\A$. Equations
\eqref{ent:eq:prep-rec1}--\eqref{ent:eq:prep-rec2} give $A=(P+R)(1+Q)$
coefficient by coefficient, and the second factor is a unit by
\cref{ent:lem:A-units}.

For uniqueness, suppose two pairs differ. The union of their positive
supports is well ordered, so there is a least exponent $\eta$ at which either
pair differs. At that exponent, equality of the products gives
\[
 P(Q_\eta-Q'_\eta)+(R_\eta-R'_\eta)=0.
\]
The second summand has degree less than $d$. Uniqueness of polynomial
division forces both differences to vanish, a contradiction. For $d=0$, the
normalization gives $P=1$, $W=1$ and $U=A$.
\end{proof}

This proof is deliberately not a claim that a Banach iteration converges; see
\cref{ent:rem:no-contraction}. It remains valid when no positive element of
$\Gamma$ has cofinal integer multiples. The uniformly bounded remainder
degree is the mechanism that extracts a finite zero scheme from an infinite
Hahn series.

\begin{corollary}[The certificate survives enlargement of the value group]
\label{ent:cor:prep-extension}
Let $\Gamma\subseteq\Delta$ be an ordered-group inclusion and regard
$A\in\A$ as an element of $\mathcal A_\Delta$ under the canonical inclusion.
The factorization $A=WU$ of \cref{ent:thm:prep} computed over $\Gamma$ is
also the factorization computed over $\Delta$: the same $W$ and $U$, with the
same support certificate $S^*\setminus\{0\}$.
\end{corollary}
\begin{proof}
The recursion \eqref{ent:eq:prep-rec1}--\eqref{ent:eq:prep-rec2} is indexed
by $S^*$, computed inside $\C[X]$, and refers to no element of the value
group beyond the exponents already occurring in $S$. It therefore produces
the same data over $\Delta$, and uniqueness over $\Delta$ identifies the two
factorizations.
\end{proof}

\Cref{ent:cor:prep-extension} is not a formality: it is exactly what lets the
same prepared polynomial be used on both sides of a scalar extension in
\cref{ent:thm:zero-conservation}, and it is available only because the
preparation was proved by a support recursion rather than by an iteration
whose convergence would have to be re-established in the larger field.

\begin{example}[A concrete preparation]\label{ent:ex:cubic}
Take $P=X^2-X$, $d=2$, and
\begin{equation}\label{ent:eq:cubic}
 A=X^2-X+t(X^3+1)\in\A,
\end{equation}
so that $S=\{1\}$ and $S^*=\N$. The recursion gives $Q_1=X+1$, $Q_2=-1$,
$Q_3=3$ and the first three remainders
\[
 R_1=X+1,\qquad R_2=-3X-1,\qquad R_3=8X+2,
\]
each of degree less than $2$ as required. For instance
$X^3+1=(X^2-X)(X+1)+(X+1)$ gives the first step, and
$H_2=-R_1Q_1=-(X+1)^2=(X^2-X)(-1)+(-3X-1)$ the second. The recursion is run
to perturbation order eight by the accompanying finite check suite
(\cref{ent:app:checks}).
\end{example}

\section{Exact root counting at every scale}
\label{ent:sec:zeros}

\subsection{Initial polynomials on a ball}
For $\gamma\in\Gamma$, define the closed valuation ball
\[
 B_\gamma=\{z\in\K:v(z)\ge\gamma\}.
\]
A larger ball corresponds to a smaller, typically negative, $\gamma$. For
$F\in\E\setminus\{0\}$, put
\begin{align}
 m_\gamma(F)&=\min_{n:a_n\ne0}\{v(a_n)+n\gamma\},\label{ent:eq:weighted}\\
 I_\gamma(F)(X)&=
 \sum_{v(a_n)+n\gamma=m_\gamma(F)}\lc(a_n)X^n.
 \label{ent:eq:initial}
\end{align}
The minimum is attained, and attained at only finitely many indices, and the
sum is a finite nonzero polynomial, precisely because the weighted
coefficient valuations diverge (\cref{ent:thm:criterion}). Equivalently
$I_\gamma(F)$ is the reduction of $t^{-m_\gamma(F)}F(t^\gamma X)$. Write
$d_\gamma(F)=\deg I_\gamma(F)$ and $e_\gamma(F)=\ord_{X=0}I_\gamma(F)$.

A zero's multiplicity is the order of the translated formal power series
$F(z+Z)\in\K[[Z]]$; the translation identity \eqref{ent:eq:taylor} makes this
intrinsic.

\begin{corollary}[Finite zero geometry on the integral disk]
\label{ent:cor:integral-zeros}
Let $A\in\A$ be nonzero with $w(A)=0$ and initial polynomial monic of degree
$d$ after normalization, and let $W$ be its prepared polynomial. The zeros of
$A$ in $\OO$ are exactly the $d$ roots of $W$ with multiplicity, all of them
integral, and for each $c\in\C$ the number of them with residue $c$ is the
multiplicity of $c$ as a root of the initial polynomial. In particular a
nonzero element of $\A$ has only finitely many zeros on the integral disk.
No complex compactness is used.
\end{corollary}
\begin{proof}
This is the content of the first two paragraphs of the proof of
\cref{ent:thm:roots} below, applied at $\gamma=0$.
\end{proof}

\begin{theorem}[Ball, shell, and residue-direction counts]
\label{ent:thm:roots}
For $F\in\E\setminus\{0\}$ and $\gamma\in\Gamma$:
\begin{enumerate}[label=\textnormal{(\alph*)},leftmargin=2em]
\item $B_\gamma$ contains exactly $d_\gamma(F)$ zeros of $F$, counted with
multiplicity.
\item The open ball $\{z:v(z)>\gamma\}$ contains exactly $e_\gamma(F)$ zeros,
and the shell $\{z:v(z)=\gamma\}$ contains $d_\gamma(F)-e_\gamma(F)$ zeros.
\item For each $c\in\C$, the zeros for which $v(z)\ge\gamma$ and
$\res(t^{-\gamma}z)=c$ have total multiplicity $\ord_{X=c}I_\gamma(F)$.
\end{enumerate}
\end{theorem}
\begin{proof}
Let $c_0$ be the leading coefficient of $I_\gamma(F)$ and apply
\cref{ent:thm:prep} to
\[
 A=\frac{F(t^\gamma X)}{c_0t^{m_\gamma(F)}}.
\]
It gives $A=WU$ with $W\in\OO[X]$ monic of degree $d_\gamma(F)$,
$\res(W)=I_\gamma(F)/c_0$, and $U$ of initial polynomial $1$. The field $\K$
is algebraically closed by the standard Hahn-field theorem
\cite[Corollary~4]{Poonen}, so $W$ has exactly its degree many roots in $\K$,
counted with multiplicity.

Every root $x$ of $W$ is integral: if $v(x)<0$, the leading term $x^d$ has
valuation $dv(x)$, strictly smaller than that of every lower-degree term
whose coefficient is integral, and such terms cannot sum to zero. On $\OO$,
evaluation of $U$ has residue $1$ and is nonzero. After translation by an
integral root, $U$ has nonzero constant term, so its formal local power
series is a unit. Hence zeros and multiplicities of $A$ in $\OO$ coincide
exactly with those of $W$.

Factor $W=\prod_j(X-x_j)$ and reduce modulo $\mm$. Then
$\res(W)=\prod_j(X-\res(x_j))$, which proves the residue-direction count.
The direction $c=0$ is exactly the open ball, and the nonzero directions make
up the shell. Rescale back by $Z=t^\gamma X$.
\end{proof}

\begin{corollary}[Units, faithfulness, and values]\label{ent:cor:units}
A nonzero entire function without zeros is a nonzero constant. Distinct
elements of $\E$ define distinct functions on $\K$. Every nonconstant element
of $\E$ takes every value in $\K$.
\end{corollary}
\begin{proof}
If a zero-free $F$ has $a_0=0$ it vanishes at zero, a contradiction. If
$a_n\ne0$ for some $n>0$, choose $\delta>0$ and
\[
 \gamma=\frac{v(a_0)-v(a_n)}{n}-\delta.
\]
Then $v(a_n)+n\gamma<v(a_0)$, so the constant coefficient is not an active
term in \eqref{ent:eq:initial} and the initial polynomial has positive
degree; \cref{ent:thm:roots} supplies a zero. Therefore $F$ is constant.

For faithfulness, a nonzero $F$ has only finitely many zeros in any
$B_\gamma$, whereas that ball contains infinitely many elements $t^\gamma c$
with $c\in\C$; apply this to the difference of two series. Finally, for any
$b\in\K$ a nonconstant $F-b$ cannot be zero-free, so $F$ takes the value $b$.
\end{proof}

\begin{lemma}[Division by a linear zero factor]\label{ent:lem:linear-division}
If $F\in\E$ and $F(r)=0$, then $F=(Z-r)G$ with $G\in\E$.
\end{lemma}
\begin{proof}
Translate to $r=0$ by \cref{ent:prop:operations}, divide the formal series by
$Z$, and translate back. Each rescaled copy $F(t^\gamma X)$ lies in $\A$, and
dividing by a monomial or by a monic linear factor with integral root
preserves $\A$ by \cref{ent:lem:A-units}; conclude by
\cref{ent:thm:criterion}(2).
\end{proof}

\subsection{Stability of the counts}
\begin{corollary}[Initial-form stability]\label{ent:cor:stability}
If $F,G\in\E$ satisfy $w((F-G)(t^\gamma X))>m_\gamma(F)$, then
$m_\gamma(G)=m_\gamma(F)$ and $I_\gamma(G)=I_\gamma(F)$. All three counts in
\cref{ent:thm:roots} are therefore unchanged.
\end{corollary}
\begin{proof}
A perturbation supported strictly above the leading exponent changes no
coefficient at that exponent; apply \cref{ent:thm:roots}.
\end{proof}

\begin{proposition}[Locally finite Newton profile]\label{ent:prop:newton}
Let $F\in\E\setminus\{0\}$ and let $[\gamma_0,\gamma_1]$ be an interval of
$\Gamma$. On that interval, $\gamma\mapsto m_\gamma(F)$ is the minimum of
\emph{finitely many} affine functions $\gamma\mapsto v(a_n)+n\gamma$. The set
of active indices changes only at the slopes
\[
 \frac{v(a_i)-v(a_j)}{j-i}\in\Gamma\qquad(i<j),
\]
and between two consecutive such values the three counts of
\cref{ent:thm:roots} are constant.
\end{proposition}
\begin{proof}
By \cref{ent:thm:criterion}(4) applied at $\gamma_0$ and at $\gamma_1$, only
finitely many indices $n$ can attain the minimum \eqref{ent:eq:weighted} for
some $\gamma$ in the interval: all but finitely many satisfy
$v(a_n)+n\gamma>m_{\gamma_0}(F)+|m_{\gamma_1}(F)-m_{\gamma_0}(F)|$ throughout
it. A minimum of finitely many affine functions is piecewise affine, and two
of them exchange dominance exactly at the displayed slope, which lies in
$\Gamma$ by divisibility. \Cref{ent:cor:stability} gives the constancy.
\end{proof}

This is a valuation statement, a non-Archimedean analogue of a Rouch\'e
stability argument; it is not a theorem about Euclidean circles in the fine
surcomplex topology. Its purpose here is to retain exact multiplicities in
the factorization argument of \cref{ent:sec:products}.

\section{Radial divisors, canonical products and factorization}
\label{ent:sec:products}

\subsection{Which zero sets are admissible?}
\begin{definition}\label{ent:def:radial}
An effective divisor on $\K$ is a function $D:\K\to\N$ with set-sized
support. It is \emph{radially finite}, equivalently \emph{ball-finite}, if
\[
 \sum_{z\in B_\gamma}D(z)<\infty\qquad(\gamma\in\Gamma).
\]
The sum means a finite total multiplicity, not a Hahn sum. Let
$\Div^+_{\mathrm{rad}}(\K)$ denote the monoid of these divisors.
\end{definition}

The two source manuscripts used the two displayed adjectives for this one
condition; we keep \emph{radially finite} and record \emph{ball-finite} as
the synonym. This is a restriction on valuation balls, not on discrete
subsets for the fine surreal topology, and the distinction is necessary:
fine discreteness alone would impose no useful control on an infinite
collection inside a fixed bounded valuation ball. Radial finiteness is
exactly what the root count of \cref{ent:thm:roots} supplies.

\begin{lemma}[Countability and escape]\label{ent:lem:divisor-countable}
An infinite radially finite divisor has countably many points counted with
multiplicity. Its nonzero points can be enumerated as $(\rho_j)_{j\ge1}$,
with repetitions for multiplicities, and
\[
 v(\rho_j)\longrightarrow-\infty \quad\hbox{coinitially in }\Gamma.
\]
In particular an infinite radially finite divisor forces
$\cf(\Gamma)=\aleph_0$; and when $\cf(\Gamma)>\aleph_0$, every radially
finite divisor is finite.
\end{lemma}
\begin{proof}
Choose countably many distinct occurrences from the infinite divisor. Their
valuations cannot have a common lower bound, for otherwise one ball would
contain infinitely many occurrences. They form a countable coinitial subset
of $\Gamma$, whose negatives are a countable cofinal subset. Balls at those
scales exhaust $\K$ and each contains only finitely many divisor
occurrences, so their union contains the whole divisor and is countable. In
any one-to-one enumeration of its occurrences, only finitely many indices
have valuation at least any fixed $\gamma$, which is the asserted escape. The
occurrence at zero, if present, has finite multiplicity and is removed before
enumerating the nonzero points.
\end{proof}

\subsection{Construction and exact factorization}
\begin{theorem}[Canonical product for a radial divisor]
\label{ent:thm:product}
Let $D\in\Div^+_{\mathrm{rad}}(\K)$, let $m=D(0)$, and enumerate its nonzero
occurrences as $\rho_j$. The product
\begin{equation}\label{ent:eq:product}
 P_D(Z)=Z^m\prod_j\left(1-\frac Z{\rho_j}\right)
\end{equation}
is a well-defined element of $\E$, independent of the enumeration. Its
divisor is exactly $D$, and the coefficient of $Z^m$ is $1$. The product is
finite when the nonzero divisor is finite. Its partial products converge to
it in the weighted Gauss valuation at every scale. \emph{No exponential
correction factors are required}: the theory is genus zero at every rank.
\end{theorem}
\begin{proof}
Only the infinite case needs a support argument. Put $b_j=\rho_j^{-1}$. Fix
$\gamma\in\Gamma$ and substitute $Z=t^\gamma X$. By
\cref{ent:lem:divisor-countable}, $v(b_j)+\gamma\to+\infty$. There are only
finitely many factors for which $v(b_j)+\gamma\le0$; set them aside. For the
rest, the supports of $b_jt^\gamma$ have a well-ordered positive union and
are locally finite, by \cref{ent:lem:cofinal-tails}.

Expand the tail product over finite subsets of indices. Every exponent
belongs to the monoid generated by that positive union. For a fixed exponent,
\cref{ent:lem:neumann} gives finitely many words in the support, and local
finiteness gives finitely many choices of factor labels carrying those
exponents. Thus only finitely many product terms contribute, and the
coefficient at that exponent is a polynomial in $X$. The tail product is an
element of $\A$ with initial polynomial $1$, hence a unit. Multiplying by the
finitely many set-aside factors and by $t^{m\gamma}X^m$ gives an element of
$\A$; the same estimate shows that the finite partial products approach it in
$m_\gamma$.

For each fixed ordinary degree, this construction sums the same finite-subset
coefficient family multiplied by $t^{k\gamma}$. Consequently the elements
obtained for different $\gamma$ are rescalings of one series in $\K[[Z]]$,
which belongs to $\E$ by \cref{ent:thm:criterion}(2). Strong regrouping also
proves independence of the enumeration.

On $B_\gamma$, the set-aside factors give precisely the roots in that ball,
while the remaining unit is nowhere zero on it. Local translation shows that
multiplicities are the displayed factor multiplicities. Since the balls cover
$\K$, the divisor is exactly $D$. Every factor of the nonzero-root product
has constant term $1$, so the coefficient of $Z^m$ is $1$.
\end{proof}

\begin{theorem}[Factorization by zeros]\label{ent:thm:factorization}
Let $F=\sum a_nZ^n\in\E\setminus\{0\}$ and let $m=\ord_{Z=0}F$. Its zero
divisor is radially finite. If $\rho_j$ are all its nonzero zeros, repeated
with their multiplicities, then
\begin{equation}\label{ent:eq:factorization}
 F(Z)=a_mZ^m\prod_j\left(1-\frac Z{\rho_j}\right).
\end{equation}
Thus the zero divisor determines $F$ up to a nonzero scalar.
\end{theorem}
\begin{proof}
Radial finiteness is \cref{ent:thm:roots}. Let $G=P_{\Div(F)}$ be the
canonical product from \cref{ent:thm:product}. Since $F$ and $G$ have the
same order $m$ at zero, the formal quotient
\[
 Q=\frac{F/Z^m}{G/Z^m}\in\K[[Z]]
\]
exists, because $G/Z^m$ has constant coefficient $1$.

We claim $Q\in\E$. Fix $\gamma$. Preparation of $F(t^\gamma X)$ and
$G(t^\gamma X)$ writes them as
\[
 F(t^\gamma X)=c_F W_F U_F,\qquad
 G(t^\gamma X)=c_G W_G U_G,
\]
with $c_F,c_G\in\K^\times$, units $U_F,U_G$ of $\A$, and monic $W_F,W_G$ all
of whose roots lie in $\OO$. Those roots are exactly the normalized roots of
$F$ and of $G$ in $B_\gamma$, with multiplicity, so $W_F=W_G$. Cancellation
in $\K[[X]]$, including the common factor $X^m$, gives
\[
 Q(t^\gamma X)=\frac{c_F}{c_G}\frac{U_F}{U_G}\in\A.
\]
This holds at every scale, so \cref{ent:thm:criterion}(2) proves the claim.

The identity $F=GQ$ is now an identity of entire functions. Orders at any
translated centre add in the formal power-series ring. Since $F$ and $G$ have
identical zero multiplicities, $Q$ has no zeros, so $Q$ is constant by
\cref{ent:cor:units}. Comparing the coefficient of $Z^m$ gives $Q=a_m$.
\end{proof}

The quotient step is important and is the reason the preparation theorem was
proved at \emph{every} scale. Equal zero sets alone do not justify dividing
two infinite series in the entire category, and no compactness or transfinite
convergence assertion replaces that argument. Note also the order of the
argument: the quotient is shown strongly evaluable at all scales
\emph{before} \cref{ent:thm:criterion} upgrades it to an entire,
valuation-convergent series. In particular the proof never assumes that a
zero-free restricted series is restricted-invertible --- by
\cref{ent:ex:ranktwo-unit} that is exactly what can fail in higher rank.

\begin{corollary}[Divisibility, GCDs, and least common multiples]
\label{ent:cor:gcd}
The map $F\mapsto\Div(F)$ induces an isomorphism of commutative monoids
\[
 (\E\setminus\{0\})/\K^\times
 \ \cong\ \Div^+_{\mathrm{rad}}(\K).
\]
For nonzero $F,G\in\E$, $G$ divides $F$ in $\E$ if and only if
$\Div(G)\le\Div(F)$ pointwise. Greatest common divisors and least common
multiples are given by the pointwise minimum and maximum of divisors. In
particular $\E$ is a GCD domain.
\end{corollary}
\begin{proof}
Surjectivity is \cref{ent:thm:product}; injectivity modulo constants is
\cref{ent:thm:factorization}. Divisors add under products by Taylor
expansion. If the divisor of $G$ is contained in that of $F$, their
difference is radially finite, so its canonical product supplies the
quotient, up to the scalar fixed by \cref{ent:thm:factorization}. Pointwise
minima and maxima of two radial divisors remain radially finite; their
canonical products therefore have the defining universal divisibility
properties of a GCD and an LCM.
\end{proof}

\begin{corollary}[Finite jets and evaluation ideals]\label{ent:cor:evaluation}
For $x\in\K$ and $r\ge1$, the kernel of the Taylor jet map
\[
 \E\longrightarrow\K[U]/(U^r),\qquad
 F\longmapsto F(x+U)\bmod U^r
\]
is $(Z-x)^r\E$, and the map is surjective. In particular evaluation at $x$
has kernel $(Z-x)\E$ and residue field $\K$.
\end{corollary}
\begin{proof}
Vanishing of the jet is exactly a zero of order at least $r$; apply
\cref{ent:cor:gcd}. Every jet is realized by a finite polynomial in $Z-x$.
\end{proof}

\begin{corollary}[Rigidity and value distribution]
\label{ent:cor:value-distribution}
Let $F\in\E$ be nonzero.
\begin{enumerate}[label=\textnormal{(\alph*)},leftmargin=2em]
\item $F$ has finitely many zeros if and only if $F$ is a polynomial; a
zero-free $F$ is a nonzero constant.
\item If $F$ is nonconstant it is surjective onto $\K$; if $F$ is
nonpolynomial, every fibre $F^{-1}(b)$ is infinite.
\item If $F$ is injective as a function on $\K$, then $F$ is affine of
degree one.
\end{enumerate}
\end{corollary}
\begin{proof}
(a) A polynomial has finitely many zeros, and conversely
\eqref{ent:eq:factorization} with a finite product is a polynomial.
Zero-freeness is \cref{ent:cor:units}. (b) Surjectivity is
\cref{ent:cor:units}; if $F$ is nonpolynomial then $F-b$ is nonpolynomial,
so by (a) it has infinitely many zeros. (c) An injective $F$ is nonconstant,
hence a polynomial by (b), say of degree $d\ge1$ with leading coefficient
$a_d$. For each $b$, injectivity makes the divisor of $F-b$ a single point
$r_b$ with multiplicity $d$, so $F(Z)-b=a_d(Z-r_b)^d$ by
\cref{ent:thm:factorization}. Comparing two values $b\ne b'$ gives
$a_d\bigl[(Z-r_b)^d-(Z-r_{b'})^d\bigr]=b'-b$, a nonzero constant. For
$d\ge2$ the left side has a nonzero coefficient of $Z^{d-1}$ unless
$r_b=r_{b'}$, and $r_b=r_{b'}$ makes it zero; either way we contradict
$b\ne b'$. Hence $d=1$.
\end{proof}

Part (b) and part (c) are familiar non-Archimedean value distribution, stated
here for completeness at arbitrary rank and not offered as independent
priority claims.

\begin{corollary}[Conjugation and real-coefficient functions]
\label{ent:cor:conjugation}
A nonzero entire series has all coefficients in $\R((t^\Gamma))$ if and only
if its zero divisor is invariant under coefficientwise conjugation and its
first nonzero coefficient lies in $\R((t^\Gamma))$.
\end{corollary}
\begin{proof}
Conjugating coefficients conjugates each value and each zero with its
multiplicity. Conversely, under the stated conditions the conjugate series
has the same divisor and the same first nonzero coefficient, so uniqueness in
\cref{ent:thm:factorization} makes the two series equal.
\end{proof}

\subsection{What ``factorization'' does not imply}
\begin{proposition}\label{ent:prop:nonnoetherian}
If $\E$ contains a nonpolynomial function, it is not Noetherian and is not
atomic. Its irreducible elements are precisely the scalar multiples of linear
polynomials.
\end{proposition}
\begin{proof}
By \cref{ent:thm:factorization}, a nonpolynomial entire function has
infinitely many zeros with multiplicity. Let $\rho_1,\rho_2,\ldots$ be an
escaping sequence of nonzero roots and put
$T_n(Z)=\prod_{j\ge n}(1-Z/\rho_j)$. Then $T_n=(1-Z/\rho_n)T_{n+1}$, so
\[
 (T_1)\subsetneq(T_2)\subsetneq(T_3)\subsetneq\cdots,
\]
strictness following either by evaluating multiplicities at $\rho_n$ or from
\cref{ent:cor:gcd}. Thus the ring is not Noetherian.

A nonunit has a zero. If its divisor has at least two occurrences, separating
one occurrence from the remaining divisor gives a product of two nonunits. A
singleton divisor gives a scalar multiple of one linear factor and is
irreducible. Consequently a function with an infinite divisor cannot be a
\emph{finite} product of irreducibles, which is the failure of atomicity.
\end{proof}

The canonical product theorem is therefore not a unique-factorization-domain
theorem. Its product is infinite in a precisely controlled Hahn sense; an
algebraic UFD requires finite irreducible factorizations.

\section{The higher-rank obstruction to interpolation and B\'ezout}
\label{ent:sec:obstruction}

\subsection{A sequence of genuinely new Archimedean scales}
For positive $\alpha,\beta\in\Gamma$, write
\[
 \beta\ggA\alpha
 \quad\Longleftrightarrow\quad
 \beta>k\alpha\text{ for every integer }k\ge1.
\]
This is Archimedean domination: a genuinely new Archimedean class, not merely
a large finite ratio.

\begin{lemma}\label{ent:lem:escaping-classes}
Suppose $\cf(\Gamma)=\aleph_0$ and $\Gamma$ has no order unit. There is a
positive cofinal sequence $(\gamma_n)_{n\ge1}$ with
$\gamma_{n+1}\ggA\gamma_n$ for every $n$.
\end{lemma}
\begin{proof}
Choose a positive cofinal sequence $(q_n)$. After selecting $\gamma_n$, the
set $\{k\gamma_n:k\ge1\}$ is not cofinal, since $\gamma_n$ is not an order
unit; hence it has a strict upper bound $b_n\in\Gamma$. Choose
$\gamma_{n+1}>\max\{b_n,q_{n+1},\gamma_n\}$ and start with
$\gamma_1\ge q_1>0$. The resulting sequence has the required domination and
is cofinal because it dominates the $q_n$.
\end{proof}

\subsection{A universal bound on values at escaping monomials}
The next lemma is the core obstruction. It is short, but it uses both halves
of \cref{ent:def:entire}: a common lower bound on the coefficient supports,
and legitimate evaluation at the large argument.

\begin{lemma}[Archimedean growth barrier]\label{ent:lem:growth-barrier}
Let $(\gamma_n)$ be positive and cofinal, and let $\lambda_n>0$ satisfy
$\lambda_n\ggA\gamma_n$ for every sufficiently large $n$. For each
$H\in\E$, eventually
\begin{equation}\label{ent:eq:growth-barrier}
 v\bigl(H(t^{-\gamma_n})\bigr)>-\lambda_n.
\end{equation}
The value $+\infty$ is allowed when $H$ vanishes at the argument.
\end{lemma}
\begin{proof}
The assertion is immediate for $H=0$. Write $H=\sum_{k\ge0}h_kZ^k$ and let
$m=\min_{h_k\ne0}v(h_k)$, which exists by strong summability at $Z=1$. For
large $n$, $\gamma_n>|m|$. Then each nonzero term at $Z=t^{-\gamma_n}$
satisfies
\[
 v(h_kt^{-k\gamma_n})\ge m-k\gamma_n
      >-(k+1)\gamma_n>-\lambda_n.
\]
The terms' valuations tend cofinally to infinity as $k\to\infty$ by
\cref{ent:thm:criterion}, and their nonempty set has a minimum attained at a
finite $k$. That minimum is strictly greater than $-\lambda_n$, and the
valuation of their Hahn sum is at least the minimum, giving
\eqref{ent:eq:growth-barrier}.
\end{proof}

\begin{corollary}[Explicit failure of value interpolation]
\label{ent:cor:no-interpolation}
For the sequence of \cref{ent:lem:escaping-classes}, the prescription
\[
 H(t^{-\gamma_n})=t^{-\gamma_{n+1}}\qquad(n\ge1)
\]
has no solution $H\in\E$. Each finite subcollection of these prescriptions
has a polynomial solution.
\end{corollary}
\begin{proof}
The valuations of the prescribed values are $-\gamma_{n+1}$, in conflict with
\cref{ent:lem:growth-barrier}. The nodes are distinct, so ordinary finite
Lagrange interpolation over the field $\K$ realizes any finite subcollection.
\end{proof}

This failure is not caused by nodes accumulating in a bounded ball: their
divisor is radially finite and is therefore the divisor of an entire function
by \cref{ent:thm:product}. Realizability of zeros is strictly weaker than
realizability of values.

\subsection{Two functions with no common zero and no B\'ezout identity}
\begin{theorem}[An explicit pointwise B\'ezout failure]
\label{ent:thm:no-bezout}
Assume the hypotheses of \cref{ent:lem:escaping-classes} and choose its
sequence $(\gamma_n)$. Define
\begin{align}
 \rho_n&=t^{-\gamma_n},&
 \sigma_n&=t^{-\gamma_n}+t^{\gamma_{n+1}},\label{ent:eq:paired-roots}\\
 f(Z)&=\prod_{n\ge1}(1-Z/\rho_n),&
 g(Z)&=\prod_{n\ge1}(1-Z/\sigma_n).\label{ent:eq:paired-products}
\end{align}
Then $f,g\in\E$, both have constant coefficient $1$, their roots are simple
and disjoint, and
\[
 \gcd(f,g)=1,\qquad (f,g)\ne\E.
\]
In particular there are no $A,B\in\E$ with $Af+Bg=1$.
\end{theorem}
\begin{proof}
We have $v(\rho_n)=v(\sigma_n)=-\gamma_n$, so both root sequences escape
every valuation ball and \cref{ent:thm:product} defines the two entire
functions. Different indices have different valuations; at the same index
$\sigma_n-\rho_n=t^{\gamma_{n+1}}\ne0$. Their zero sets are consequently
disjoint and simple, and \cref{ent:cor:gcd} gives their GCD.

We compute $v(g(\rho_n))$ exactly. For $j<n$, $\rho_n/\sigma_j$ has negative
valuation, so
\[
 v(1-\rho_n/\sigma_j)=\gamma_j-\gamma_n.
\]
For $j=n$,
\[
 1-\rho_n/\sigma_n=\frac{\sigma_n-\rho_n}{\sigma_n},\qquad
 v(1-\rho_n/\sigma_n)=\gamma_{n+1}+\gamma_n.
\]
For $j>n$ the ratio has positive valuation, and the infinite tail product has
residue $1$ and valuation zero by the positive-support product argument used
in \cref{ent:thm:product}. Therefore
\begin{equation}\label{ent:eq:lambda}
 \lambda_n:=v(g(\rho_n))
   =\gamma_{n+1}+(2-n)\gamma_n+\sum_{j<n}\gamma_j.
\end{equation}
The sum here is finite. Since $0<\gamma_j<\gamma_n$ for $j<n$, the part of
\eqref{ent:eq:lambda} other than $\gamma_{n+1}$ is bounded in absolute value
by a finite multiple of $\gamma_n$. Hence $\lambda_n\ggA\gamma_n$.

Suppose $Af+Bg=1$. Evaluate at $\rho_n$. The values are well defined and
evaluation is multiplicative, so
\[
 B(\rho_n)=\frac1{g(\rho_n)},\qquad
 v(B(\rho_n))=-\lambda_n.
\]
For all sufficiently large $n$ this contradicts
\cref{ent:lem:growth-barrier}. Thus the ideal is proper.
\end{proof}

The cancellation-sensitive step is \eqref{ent:eq:lambda}. No unproved
assertion that nearby zeros force a small value is used: the contribution of
each earlier root is retained, and the infinite tail is proved to be a unit.
The earlier-root terms can lower the valuation by finite multiples of
$\gamma_n$, but cannot offset the new Archimedean scale $\gamma_{n+1}$.

\begin{corollary}[A finitely generated ideal invisible to point tests]
\label{ent:cor:invisible-ideal}
In the situation of \cref{ent:thm:no-bezout} there is a proper two-generated
ideal contained in no evaluation maximal ideal $(Z-x)\E$. Every maximal ideal
containing it is therefore not the kernel of evaluation at any $x\in\K$.
\end{corollary}
\begin{proof}
The ideal $(f,g)$ is proper but $f$ and $g$ do not vanish simultaneously. By
\cref{ent:cor:evaluation}, containment in $(Z-x)\E$ would mean that both
vanish at $x$. A maximal ideal containing a proper ideal exists by Zorn's
lemma; none of these maximal ideals is an evaluation ideal.
\end{proof}

This is a failure of the assertion that \emph{every finitely generated proper
ideal has a common zero}. We do not call it a metric corona counterexample,
because no norm lower bound and no bounded-function algebra has been
specified. Nor do we claim that non-evaluation maximal ideals are absent in
the rank-one entire ring; the statement concerns a \emph{two-generated} ideal
already lacking a common zero.

\subsection{The missing image of the infinite jet map}
\label{ent:subsec:jet-image}
Let $D$ be the simple divisor with nodes $\rho_n$. The evaluation map
\[
 \operatorname{ev}_D:\E\longrightarrow\prod_{n\ge1}\K,
 \qquad H\longmapsto(H(\rho_n))_n
\]
has kernel $f\E$ by \cref{ent:cor:gcd}, and it is not surjective by
\cref{ent:cor:no-interpolation}. In particular
\[
 \E/(f)\hookrightarrow\prod_{n\ge1}\K
\]
is a proper subring. The element $g\bmod f$ maps to an everywhere nonzero
sequence, hence to a unit of the product ring, but it is not a unit in
$\E/(f)$: an inverse there would give representatives satisfying a B\'ezout
identity.

Thus the elementary implication ``invertible at each node implies invertible
in the global quotient'' is exactly where unrestricted infinite Chinese
remainder reasoning fails. Every finite projection of this map is surjective.
Neither a finite truncation nor a finite root calculation can detect the
global failure; see \cref{ent:app:checks}.

\section{The positive case: an order unit restores interpolation}
\label{ent:sec:positive}

\subsection{A rank-one coarsening with the same entire functions}
It would be incorrect to infer that the preceding obstruction occurs for
every higher-rank group. A group may have several convex levels and still
have a greatest Archimedean class. In that case a single cofinal scale
permits ordinary non-Archimedean analytic methods.

\begin{lemma}[Real-valued coarsening]\label{ent:lem:coarsening}
Suppose $h>0$ is an order unit of $\Gamma$. There is an order-preserving
homomorphism $\pi:\Gamma\to\R$ with $\pi(h)=1$, whose kernel is the convex
subgroup
\[
 H=\{\gamma\in\Gamma:n|\gamma|<h\text{ for all }n\ge1\}.
\]
The formula
\[
 |a|_\pi=\exp(-\pi(v(a)))\quad(a\ne0),\qquad |0|_\pi=0
\]
defines a nontrivial real-valued non-Archimedean absolute value on $\K$. Its
topology agrees with the intrinsic valuation topology, and $\K$ is complete
for it. Furthermore
\begin{equation}\label{ent:eq:entire-coarsening}
 \E=\left\{\sum_{n\ge0}a_nZ^n:
           |a_n|_\pi R^n\longrightarrow0\text{ for every real }R>0
     \right\}.
\end{equation}
\end{lemma}
\begin{proof}
Being an order unit bounds every $\gamma$ above and below by integer
multiples of $h$. Define
\[
 \pi(\gamma)=\sup\{q\in\Q:qh\le\gamma\}\in\R.
\]
Rational upper and lower approximations show
$\pi(\alpha+\beta)=\pi(\alpha)+\pi(\beta)$ and $\pi(h)=1$; monotonicity is
immediate. The kernel consists exactly of the elements infinitesimal relative
to $h$, giving the displayed $H$. In particular $\pi(\gamma)>n$ implies
$\gamma>nh$, and $\gamma>(n+1)h$ implies $\pi(\gamma)>n$. The balls specified
by integer multiples of $h$ are cofinal among all valuation neighbourhoods of
zero, and these inequalities show that the real-valued and ordered-group
topologies agree.

For completeness, take a Cauchy sequence for $|\cdot|_\pi$ and choose a
subsequence $(x_{j_n})$ with $v(x_{j_{n+1}}-x_{j_n})>nh$. The differences
have valuations tending cofinally to infinity, so their Hahn sum exists by
\cref{ent:lem:cofinal-tails}; adding it to $x_{j_1}$ gives a limit of the
subsequence, and the Cauchy property gives the same limit for the whole
sequence. Algebraic closedness of $\K$ was already supplied by the Hahn-field
theorem.

By \cref{ent:thm:criterion}, the left side of
\eqref{ent:eq:entire-coarsening} is characterized by $v(a_n)/n\to+\infty$ in
$\Gamma$. This is equivalent to $\pi(v(a_n))/n\to+\infty$ in $\R$: integer
multiples of $h$ test the forward implication, and for the reverse one
chooses a real bound strictly exceeding $\pi(\gamma)$ for a prescribed
$\gamma\in\Gamma$. The usual elementary coefficient test identifies this last
condition with the right side of \eqref{ent:eq:entire-coarsening}:
explicitly, it implies $\pi(v(a_n))-n\log R\to+\infty$ for every $R$;
conversely, failure of the ratio condition gives infinitely many ratios at
most some $C$, and then $R=e^{C+1}$ violates the coefficient test.
\end{proof}

This is a \emph{coarsening}, not a valuation-preserving embedding of $\Gamma$
in $\R$. The kernel $H$ may be nonzero, and residue-direction data is not
identified across the two valuations: the coarsening changes residue
information even while preserving the topology and the entire ring. It is the
cofinality of the topology and the all-radius coefficient test that permit
its use here.

\subsection{A complete Hermite interpolation argument}
The next proof records the damping mechanism that fails in
\cref{ent:sec:obstruction}. It is included instead of importing an
unqualified analytic interpolation theorem.

\begin{theorem}[Arbitrary radial Hermite interpolation]
\label{ent:thm:interpolation}
Assume $\Gamma$ has an order unit. Let $(x_n)$ be distinct points of a
radially finite set, and assign a positive integer $m_n$ to each point so
that the divisor $\sum_nm_n[x_n]$ is radially finite. For arbitrary jets
$J_n\in\K[U]/(U^{m_n})$ there exists $F\in\E$ with
\[
 F(x_n+U)\equiv J_n(U)\pmod{U^{m_n}}
 \qquad\text{for every }n.
\]
\end{theorem}
\begin{proof}
Use $|\cdot|=|\cdot|_\pi$ from \cref{ent:lem:coarsening}. A radially finite
infinite set can be enumerated with $|x_n|\to\infty$: a ball of any fixed
real absolute-value radius is contained in some original valuation ball, and
conversely an original ball has bounded $|\cdot|_\pi$ radius. Handle the
finitely many nodes with $|x_n|\le1$, including zero, first by a finite
Hermite polynomial, and include their factors in every later vanishing
product. Reindex the remaining nodes by $n\ge1$, so that $|x_n|>1$ and
$|x_n|\to\infty$.

Suppose a polynomial $S_{n-1}$ has the required jets at earlier nodes and put
\[
 Q_{n-1}(Z)=\prod_{j<n}(Z-x_j)^{m_j},
\]
including the finitely many initially handled nodes. It is a unit in the
local jet ring at $x_n$, so there is a unique jet
\[
 B_n(U)=\frac{J_n(U)-S_{n-1}(x_n+U)}{Q_{n-1}(x_n+U)}
              \pmod{U^{m_n}}.
\]
For a positive integer $N$ define the polynomial of degree less than $m_n$
\begin{equation}\label{ent:eq:interpolation-r}
 r_{n,N}(U)=B_n(U)(1+U/x_n)^{-N}\pmod{U^{m_n}}
\end{equation}
and the polynomial correction
\begin{equation}\label{ent:eq:interpolation-p}
 p_{n,N}(Z)=Q_{n-1}(Z)(Z/x_n)^N r_{n,N}(Z-x_n).
\end{equation}
It vanishes to the required orders at all earlier nodes, and by
\eqref{ent:eq:interpolation-r} its jet at $x_n$ is the required residual. The
parameter $N$ can therefore be used purely to control its size.

Set $R_n=|x_n|^{1/2}$, so $1<R_n<|x_n|$ and $R_n\to\infty$. For a polynomial
$P=\sum p_kZ^k$ write $\|P\|_R=\max_k|p_k|R^k$. If
$B_n(U)=\sum_{j<m_n}b_jU^j$, the coefficient of $U^\ell$ in
\eqref{ent:eq:interpolation-r} has absolute value at most
$\max_{j\le\ell}|b_j|\,|x_n|^{j-\ell}$. This bound is independent of $N$,
because the binomial coefficients $\binom{-N}{k}$ are integers and so have
absolute value at most $1$. Since $R_n<|x_n|$, substitution $U=Z-x_n$ gives
\[
 \|r_{n,N}(Z-x_n)\|_{R_n}
 \le\max_{j<m_n}|b_j|\,|x_n|^j=:C'_n,
\]
and consequently
\begin{equation}\label{ent:eq:damping}
 \|p_{n,N}\|_{R_n}
 \le \|Q_{n-1}\|_{R_n} C'_n
          \left(\frac{R_n}{|x_n|}\right)^N.
\end{equation}
The right side tends to zero as the ordinary integer $N$ tends to infinity.
Choose $N_n$ so that this norm is at most $2^{-n}$ and put
$S_n=S_{n-1}+p_{n,N_n}$.

For every fixed real $R>0$, eventually $R\le R_n$, so
$\|p_{n,N_n}\|_R\le2^{-n}$. Hence $(S_n)$ is Cauchy in the Gauss norm at
every radius. The space of series satisfying $|a_k|R^k\to0$ is complete for
that norm: coefficientwise limits exist by completeness of $\K$, and uniform
approximation by a finite polynomial bounds the tail of the limit. Thus the
limits for all radii are one common coefficient series $F$ satisfying
\eqref{ent:eq:entire-coarsening}, so $F\in\E$.

Finally, every specified finite jet is preserved. For a fixed node $x$ and
order $m$, choose $R>|x|$. Coefficients of the translated polynomial up to
degree $m-1$ are continuous for $\|\cdot\|_R$, by the finite binomial formula
and its convergent tail. All sufficiently late $S_n$ have the prescribed jet,
so their limit does too.
\end{proof}

\begin{theorem}[B\'ezout in the order-unit case]\label{ent:thm:bezout}
If $\Gamma$ has an order unit then $\E$ is a B\'ezout domain. In particular,
if $f,g\in\E$ have no common zero there are $A,B\in\E$ with $Af+Bg=1$.
\end{theorem}
\begin{proof}
First suppose $f,g$ are nonzero with no common zero. At each zero $x$ of $f$
let $m_x$ be its multiplicity. Since $g(x)\ne0$, the formal inverse of
$g(x+U)$ exists modulo $U^{m_x}$. By \cref{ent:thm:interpolation} choose
$B\in\E$ whose jet is that inverse at every such $x$. Then $1-Bg$ has a zero
of order at least $m_x$ at each zero of $f$, so it is divisible by $f$ in
$\E$ by \cref{ent:cor:gcd}; write $1-Bg=Af$.

For arbitrary nonzero $f,g$, divide both by their GCD $d$, which exists by
\cref{ent:cor:gcd}. The quotients have no common zero, so their ideal is the
unit ideal; multiplying the resulting identity by $d$ shows $(f,g)=(d)$.
Cases involving zero are immediate. Induct on the finite number of generators
to obtain the B\'ezout property.
\end{proof}

For a radially finite divisor $D$ the full jet map is therefore surjective:
\[
 \E\big/\bigl(P_D\bigr)
 \ \cong\ \prod_{x\in\supp D}\K[U]/(U^{D(x)}).
\]
The kernel is again given by divisibility, and surjectivity is
\cref{ent:thm:interpolation}. This exact unrestricted Chinese remainder
statement contrasts sharply with the proper image of
\cref{ent:subsec:jet-image}.

\section{Change of workspace: the universal scalar-extension domain}
\label{ent:sec:extension}

The adjective ``entire'' must name its field of arguments. A formal series
with coefficients in one Hahn field may be entire there and fail even strong
summability at a single point of a larger field. This section gives the exact
boundary, and then the exact description of everything inside it.

Throughout, $\Gamma\subseteq\Delta$ are nonzero divisible ordered abelian
groups with the canonical inclusion $\K\subseteq\KD$, both full Hahn fields
over the \emph{same} coefficient field $\C$.

\subsection{The sharp criterion for remaining entire}
\begin{theorem}[Sharp change-of-workspace criterion]
\label{ent:thm:extension-criterion}
With the above notation,
\begin{equation}\label{ent:eq:extension}
 \mathcal E_\Delta\cap\K[[Z]]=
 \begin{cases}
   \E,&\Gamma\text{ is cofinal in }\Delta,\\
   \K[Z],&\Gamma\text{ is not cofinal in }\Delta.
 \end{cases}
\end{equation}
In the second case a single monomial argument of $\KD$ makes \emph{every}
nonpolynomial series with coefficients in $\K$ fail strong summability.
\end{theorem}
\begin{proof}
If $\Gamma$ is cofinal in $\Delta$, a sequence of elements of
$\Gamma\cup\{+\infty\}$ tends cofinally to infinity in $\Gamma$ if and only
if it does so in $\Delta$. Apply \cref{ent:thm:criterion} to the coefficient
slopes.

If $\Gamma$ is not cofinal, choose $\delta\in\Delta$ strictly larger than
every element of $\Gamma$. For a nonpolynomial $F=\sum a_nZ^n\in\K[[Z]]$
consider evaluation at $Z=t^{-\delta}$. For any two nonzero coefficients with
$m>n$,
\[
 (v(a_m)-m\delta)-(v(a_n)-n\delta)
       =v(a_m)-v(a_n)-(m-n)\delta<0,
\]
since $v(a_m)-v(a_n)\in\Gamma$ and $(m-n)\delta\ge\delta$. The leading
exponents of the nonzero terms therefore form an infinite strictly descending
sequence, so their support union is not well ordered. No nonpolynomial series
over $\K$ can be strongly entire over $\KD$. Polynomials are entire over
either field, proving \eqref{ent:eq:extension}.
\end{proof}

\subsection{The convex-hull valuation ring}
\begin{definition}\label{ent:def:hull}
The convex hull of $\Gamma$ in $\Delta$ is
\begin{equation}\label{ent:eq:convex-hull}
 H=\conv_\Delta(\Gamma)
 =\{s\in\Delta:|s|\le\gamma\text{ for some }\gamma\in\Gamma_{>0}\},
\end{equation}
a divisible convex subgroup of $\Delta$. Coarsening $v_\Delta$ by $H$ gives a
valuation with value group $\Delta/H$.
\end{definition}

\begin{lemma}\label{ent:lem:coarsened-ring}
The valuation ring of that coarsening is
\begin{equation}\label{ent:eq:D}
 \Dom
 =\{0\}\cup\{z\ne0:v_\Delta(z)\ge\gamma
                       \text{ for some }\gamma\in\Gamma\}.
\end{equation}
Its maximal ideal consists of zero together with the elements satisfying
$v_\Delta(z)>\gamma$ for every $\gamma\in\Gamma$. Its residue field is
canonically $\C((t^H))$. Moreover $\Dom=\KD$ exactly when $\Gamma$ is cofinal
in $\Delta$.
\end{lemma}
\begin{proof}
A value $s$ has nonnegative class modulo $H$ exactly when it is bounded below
by some member of $H$, and every member of $H$ is itself bounded below by a
member of $\Gamma$; this proves \eqref{ent:eq:D}. The class is strictly
positive exactly when $s>H$, equivalently $s>\Gamma$.

For an element of this valuation ring, discard all support exponents outside
$H$. No negative coset can occur, and the discarded positive cosets form the
coarsened maximal ideal. The remaining Hahn series has support in $H$. This
is a surjective residue homomorphism onto $\C((t^H))$ with the stated kernel.
Finally the ring is the whole field exactly when $H=\Delta$, which by
\eqref{ent:eq:convex-hull} is equivalent to cofinality of $\Gamma$ in
$\Delta$.
\end{proof}

The residue field here is generally \emph{not} $\C$. It records all scales in
the convex hull $H$, whereas the original fine valuation's residue field
records only the exponent zero.

\subsection{Exact strong evaluation, not merely a sufficient domain}
\begin{theorem}[Universal scalar-extension domain]\label{ent:thm:extension}
Let $f=\sum a_nZ^n\in\E$ be nonpolynomial. In $\KD$ the family $(a_nz^n)_n$
is strongly summable if and only if $z\in\Dom$. The domain is therefore
independent of the choice of nonpolynomial $f$, and the values on it again
lie in $\Dom$.
\end{theorem}
\begin{proof}
Suppose $z\ne0$ lies in the stated ring. Choose $\gamma\in\Gamma$ with
$v_\Delta(z)\ge\gamma$ and put $u=t^{-\gamma}z\in\mathcal O_\Delta$. Since
$f$ is entire over $\K$, \cref{ent:thm:criterion}(2) gives
$f(t^\gamma X)\in\A\subseteq\mathcal A_\Delta$, and
\cref{ent:lem:A-units} gives strong summability of its evaluation family at
the integral point $u$. That family is the original family $(a_nz^n)_n$. With
$\mu=m_\gamma(f)\in\Gamma$, the same lemma gives $v_\Delta(f(z))\ge\mu$
unless the value is zero, so the value lies in $\Dom$. At $z=0$ the assertion
is immediate.

Conversely suppose $s=v_\Delta(z)<\gamma$ for every $\gamma\in\Gamma$. List
the infinitely many nonzero coefficients at increasing indices
$n_1<n_2<\cdots$ and write $\alpha_j=v(a_{n_j})\in\Gamma$. Since $s<0$,
\begin{align*}
 (\alpha_{j+1}+n_{j+1}s)-(\alpha_j+n_js)
 &=\alpha_{j+1}-\alpha_j+(n_{j+1}-n_j)s\\
 &\le\alpha_{j+1}-\alpha_j+s<0,
\end{align*}
the final inequality because $s$ is below the element
$\alpha_j-\alpha_{j+1}$ of $\Gamma$. These are actual leading exponents of
the nonzero summands $a_{n_j}z^{n_j}$, regardless of the higher terms of $z$.
Their strict descent prevents the support union from being well ordered, so
the family is not strongly summable.
\end{proof}

\begin{corollary}[Cofinal extensions are exactly the entire extensions]
\label{ent:cor:entire-extension}
For nonpolynomial $f\in\E$, the same coefficient series belongs to
$\mathcal E_\Delta$ if and only if $\Gamma$ is cofinal in $\Delta$.
Polynomials always survive.
\end{corollary}
\begin{proof}
Combine \cref{ent:thm:extension} with \cref{ent:lem:coarsened-ring};
polynomial evaluations are finite sums.
\end{proof}

\begin{remark}[The negative and the positive statement are compatible]
\label{ent:rem:reconcile}
\Cref{ent:thm:extension-criterion} and \cref{ent:thm:extension} must be read
together, and \cref{ent:rem:conflict} already stated how. The first says that
in the noncofinal case a nonpolynomial $f$ is \emph{not entire over} $\KD$,
so that $\mathcal E_\Delta\cap\K[[Z]]$ collapses to $\K[Z]$; the second says
that the very same series remains strongly summable on the proper subring
$\Dom$, with values in $\Dom$, and \cref{ent:thm:zero-conservation} below
says it keeps every one of its zeros there. There is no disagreement about
whether the function survives: it survives on $\Dom$ and nowhere else, and
``not entire over $\KD$'' is precisely the statement $\Dom\ne\KD$. Indeed the
noncofinal half of \cref{ent:thm:extension-criterion} is the special case
$z=t^{-\delta}$ of the necessity half of \cref{ent:thm:extension}, and its
independent descending-exponent proof is retained because it exhibits the
single witnessing point explicitly.
\end{remark}

A cofinal extension need not preserve rank, and a noncofinal one need not
introduce any algebraic element: the obstruction lies entirely in the ordered
hierarchy of scales. Nor does a strongly defined value in the larger field
have to be a limit there. If $\Delta$ contains a positive value above all of
$\Gamma$, the original cofinal coefficient valuations in $\Gamma$ no longer
tend to infinity in $\Delta$.

\subsection{No new zeros, even for an infinite divisor}
\begin{theorem}[Zero conservation under arbitrary enlargement]
\label{ent:thm:zero-conservation}
Let $0\ne f\in\E$. Every zero of its strongly defined extension on $\Dom$
belongs to $\K$, with unchanged multiplicity. This remains true when $f$ has
infinitely many zeros.
\end{theorem}
\begin{proof}
Suppose $z\in\Dom$ and $f(z)=0$. Choose $\gamma\in\Gamma$ with
$v_\Delta(z)\ge\gamma$. Preparation of $f(t^\gamma X)$ over $\K$
(\cref{ent:thm:prep}) gives
\[
 f(t^\gamma X)=c\,t^{\mu}W_\gamma(X)U_\gamma(X),
\]
with $W_\gamma$ monic over $\OO$ and $U_\gamma$ a unit of $\A$ with initial
polynomial $1$. By \cref{ent:cor:prep-extension} this is also the preparation
over $\Delta$, and by \cref{ent:lem:A-units} the unit factor is nonzero at
every point of $\mathcal O_\Delta$. Therefore $W_\gamma(t^{-\gamma}z)=0$. The
polynomial $W_\gamma$ already splits into linear factors over the
algebraically closed field $\K$, so $t^{-\gamma}z$, and hence $z$, lies in
$\K$.

The same prepared polynomial has the same linear-factor multiplicities in
both fields, and the evaluated unit is nonzero in both, so local orders of
vanishing are unchanged. The argument uses only the finitely many zeros in
one containing ball and applies separately to each zero; it does not require
the global divisor to be finite.
\end{proof}

This is the one statement of the section that imports algebraic closedness.
\Cref{ent:app:audit} records that dependency.

\subsection{No alternative entire continuation}
\begin{proposition}[Identity on constants]\label{ent:prop:identity}
Two formal series over $\KD$ lying in $\mathcal A_\Delta$ and agreeing at
infinitely many distinct points of $\mathcal O_\Delta$ are equal as formal
series. In particular agreement at all ordinary complex constants suffices.
\end{proposition}
\begin{proof}
Their difference lies in $\mathcal A_\Delta$. If it were nonzero,
\cref{ent:cor:integral-zeros} would give it only finitely many zeros in the
integral disk.
\end{proof}

\begin{theorem}[Obstruction to any entire-series continuation]
\label{ent:thm:no-continuation}
Suppose $\Gamma$ is not cofinal in $\Delta$ and $f\in\E$ is nonpolynomial.
There is no $g\in\mathcal E_\Delta$ agreeing with $f$ at all points of $\K$.
In fact there is no such $g$ agreeing with $f$ merely at all points of $\C$.
\end{theorem}
\begin{proof}
The coefficient series of $f$ belongs to $\mathcal A_\Delta$, since its
strong support certificate at scale zero over $\Gamma$ remains valid under an
ordered embedding; the same is true of $g$. Agreement on $\C$ would force the
two coefficient series to be identical by \cref{ent:prop:identity}. But
\cref{ent:cor:entire-extension} says that this coefficient series is not
entire over $\KD$, a contradiction.
\end{proof}

This rules out an alternative entire \emph{power-series} continuation. It
does not rule out arbitrary functions extending the pointwise values, or
continuations in a different analytic category, and it asserts no density of
$\C$ in any fine surreal topology.

\subsection{The boundary is intrinsic to the function}
\begin{corollary}[Intrinsic description of the extension boundary]
\label{ent:cor:intrinsic-boundary}
Let $f\in\E$ be nonpolynomial. In any larger ordered group $\Delta$ the
following three subsets generate the same convex subgroup:
\[
 \Gamma,\qquad
 \{v(a_n)/n:n\ge1,\ a_n\ne0\},\qquad
 \{-v(r):r\ne0,\ f(r)=0\}.
\]
Consequently the universal domain of \cref{ent:thm:extension} can be
recovered either from the coefficient slopes alone or from the zero divisor
alone.
\end{corollary}
\begin{proof}
All the displayed values lie in $\Gamma$. The slopes are cofinal in $\Gamma$
by \cref{ent:thm:criterion}. There are infinitely many zeros by
\cref{ent:cor:value-distribution}(a), and their negated valuations are
cofinal by \cref{ent:lem:divisor-countable}. A cofinal subset of an ordered
group generates its full convex hull in any containing ordered group.
\end{proof}

This makes the extension obstruction intrinsic. It is not an artefact of
having embedded the coefficients in an unnecessarily large workspace whose
extra scales the function never detects.

\subsection{Consequences for the B\'ezout counterexample and for \texorpdfstring{$\No[i]$}{No[i]}}
\begin{corollary}[No repair of the counterexample by a Hahn extension]
\label{ent:cor:no-repair}
Let $f,g$ be the pair of \cref{ent:thm:no-bezout}. If $\Gamma$ is cofinal in
$\Delta$ then $f,g\in\mathcal E_\Delta$ and still no
$A,B\in\mathcal E_\Delta$ satisfy $Af+Bg=1$. If $\Gamma$ is not cofinal in
$\Delta$ then neither $f$ nor $g$ is entire on $\KD$, although both remain
strongly defined on $\mathcal D_{\Gamma,\Delta}$ with all their zeros
(\cref{ent:thm:extension,ent:thm:zero-conservation}).
\end{corollary}
\begin{proof}
In the cofinal case $(\gamma_n)$ remains cofinal in $\Delta$ and
$\gamma_{n+1}\ggA\gamma_n$ is unchanged, so the evaluation formula
\eqref{ent:eq:lambda} is unchanged. The proof of
\cref{ent:lem:growth-barrier} now applies to any candidate
$B\in\mathcal E_\Delta$, even with coefficients outside $\K$, giving the same
contradiction. In the noncofinal case apply
\cref{ent:thm:extension-criterion} to the two nonpolynomial series, and
\cref{ent:thm:extension,ent:thm:zero-conservation} for the positive half.
\end{proof}

\begin{corollary}[Whole-class rigidity on the surcomplex numbers]
\label{ent:cor:all-no}
Let $F(Z)=\sum_{n\ge0}a_nZ^n$ be an ordinary power series with coefficients
in $\No[i]$. If the family $(a_nz^n)_n$ is strongly normal-form-summable for
every $z\in\No[i]$, then only finitely many $a_n$ are nonzero.
\end{corollary}
\begin{proof}
The union of the normal-form exponent supports of the real and imaginary
parts of the countably many coefficients is a \emph{set}. Its divisible
additive span together with $1$, with signs chosen to match
$t^\gamma=\omega^{-\gamma}$, is a nonzero set-sized ordered subgroup
$\Gamma\subseteq\No$, and all $a_n$ lie in $\K$ under
\eqref{ent:eq:embedding}; the hypothesis makes $F$ entire over $\K$. Every
set of surreal numbers has a surreal strict upper bound: the set-sized cut
$\delta=\{\,\Gamma\cup\{0\}\mid\varnothing\,\}$ supplies a positive
$\delta\in\No$ larger than every member of $\Gamma$. Enlarge $\Gamma$ to its
divisible span with $\delta$. If $F$ were nonpolynomial, the
descending-leading-exponent argument of
\cref{ent:thm:extension-criterion} --- equivalently
\cref{ent:thm:extension} --- would prohibit strong evaluation at
$z=\omega^{\delta}=t^{-\delta}$, contradicting the hypothesis.
\end{proof}

\begin{remark}[This rigidity is a re-derivation, not a new theorem]
\label{ent:rem:rigidity-credit}
\Cref{ent:cor:all-no} is presented as a consequence of the extension theorem,
not as an independent contribution. The same conclusion is already in this
collection twice over: the \emph{Analysis} report proves all-scale polynomial
rigidity directly, and the \emph{Foundations} report records the exact
boundary of the phenomenon with a pair of series over $\C((t^\Q))$ differing
only in the sign of an exponent. Both source manuscripts merged here credit
the repository's pre-existing all-scale theme rather than claiming the
corollary. The step that makes the reduction legitimate is the localization
principle recorded in the foundations report --- every \emph{set} of
surcomplex numbers lies in one divisible set-sized workspace --- and the
reduction consumes exactly that principle plus the cut $\delta$.
What is proposed as new here is only the exact cofinal/noncofinal criterion
for a specified pair of set-sized workspaces
(\cref{ent:thm:extension-criterion,ent:thm:extension}) and the persistence of
the B\'ezout obstruction under every cofinal extension
(\cref{ent:cor:no-repair}). See \cref{ent:sec:collection}.
\end{remark}

\section{Assembly of the classifications}
\label{ent:sec:classification}

\begin{proof}[Proof of \cref{ent:thm:main}]
The ring is an integral domain because it is a subring of $\K[[Z]]$. Its
units are exactly the nonzero constants by \cref{ent:cor:units}: a unit is
zero-free, and nonzero constants are invertible. GCDs exist by
\cref{ent:cor:gcd}.

If $\Gamma$ has no countable cofinal subset, \cref{ent:cor:cofinality} gives
$\E=\K[Z]$, which is a PID over a field. If $\Gamma$ has an order unit, its
positive integer multiples form a countable cofinal subset, so nonpolynomial
functions exist, and \cref{ent:thm:bezout} supplies the B\'ezout property. If
$\Gamma$ has countable cofinality and no order unit,
\cref{ent:thm:no-bezout} supplies a proper ideal generated by two coprime
functions. Such an ideal cannot be principal, since any principal generator
would be a common divisor and hence a unit. Non-Noetherianity in both
nonpolynomial cases is \cref{ent:prop:nonnoetherian}.

These cases are mutually exclusive and exhaustive. For $\Gamma=0$, strong
summability at $1$ says that only finitely many coefficient supports may
contain $0$, hence only finitely many coefficients are nonzero.
\end{proof}

\begin{proof}[Proof of \cref{ent:thm:main-extension}]
Part (a) is \cref{ent:thm:extension} together with
\cref{ent:lem:coarsened-ring}; part (b) is
\cref{ent:cor:entire-extension}, whose coarse form is
\cref{ent:thm:extension-criterion}; part (c) is
\cref{ent:thm:zero-conservation}; part (d) is
\cref{ent:thm:no-continuation}; part (e) is
\cref{ent:cor:intrinsic-boundary}. The residue field statement is
\cref{ent:lem:coarsened-ring}.
\end{proof}

\begin{center}
\small
\begin{tabular}{@{}>{\raggedright\arraybackslash}p{0.34\textwidth}>{\raggedright\arraybackslash}p{0.23\textwidth}>{\raggedright\arraybackslash}p{0.34\textwidth}@{}}
\toprule
Value-group condition & Entire series & Ideal structure\\
\midrule
No countable cofinal subset & Only polynomials & PID\\[3pt]
An order unit exists & Nonpolynomial series exist &
B\'ezout, non-Noetherian\\[3pt]
Countable cofinality, no order unit & Nonpolynomial series exist &
GCD, not B\'ezout; two-generated proper ideal with no common zero\\
\bottomrule
\end{tabular}
\end{center}

\section{Explicit higher-rank and surreal examples}
\label{ent:sec:examples}

\subsection{An explicit rank-infinite subgroup of the surreal field}
\label{ent:subsec:gamma-infinity}
Let
\begin{equation}\label{ent:eq:concrete-group}
 \Gamma_\infty=\bigoplus_{j\ge0}\Q e_j,
 \qquad e_j=\omega^j\in\No,
\end{equation}
with the order inherited from $\No$: every element has finite support, and
the largest index with a nonzero coefficient determines its sign. Thus
$e_{j+1}\ggA e_j$ and $(e_j)$ is cofinal, so $\cf(\Gamma_\infty)=\aleph_0$.
No element is an order unit: the largest index in its finite support can
always be exceeded. This is case (iii) of \cref{ent:thm:main}, and by
\cref{ent:cor:unit-structure} also
$\mathcal T_{\Gamma_\infty}^\times=K_{\Gamma_\infty}^\times$.

\subsubsection*{The B\'ezout counterexample realized}
Take $\gamma_n=e_n$ for $n\ge1$, a positive cofinal sequence with
$\gamma_{n+1}\ggA\gamma_n$, as \cref{ent:lem:escaping-classes} demands.
Using $t^\gamma=\omega^{-\gamma}$, \cref{ent:thm:no-bezout} gives roots
\begin{equation}\label{ent:eq:surreal-roots}
 \rho_n=\omega^{\omega^n},\qquad
 \sigma_n=\omega^{\omega^n}+\omega^{-\omega^{n+1}}
 \qquad(n\ge1),
\end{equation}
and the functions
\begin{align}
 f(Z)&=\prod_{n\ge1}\bigl(1-\omega^{-\omega^n}Z\bigr)
      =\prod_{n\ge1}\bigl(1-t^{e_n}Z\bigr),
      \label{ent:eq:f-explicit}\\
 g(Z)&=\prod_{n\ge1}\left(
 1-\frac{Z}{\omega^{\omega^n}+\omega^{-\omega^{n+1}}}\right)
\end{align}
are entire on the fixed surcomplex Hahn field $K_{\Gamma_\infty}$, not on all
of $\No[i]$. They have real Hahn coefficients and all the displayed roots are
real surreal numbers; by \cref{ent:cor:conjugation} this is visible from the
conjugation-invariance of their divisors. The failure of a B\'ezout identity
remains true even when complex Hahn coefficients are allowed in $A,B$, and
consequently also when $A,B$ are restricted to real Hahn coefficients.

The valuation formula \eqref{ent:eq:lambda} becomes
\[
 v(g(\rho_n))=e_{n+1}+(2-n)e_n+\sum_{1\le j<n}e_j,
\]
whose first values are
\begin{align*}
 \lambda_1&=e_2+e_1, & \lambda_2&=e_3+e_1,\\
 \lambda_3&=e_4-e_3+e_2+e_1,
 &\lambda_4&=e_5-2e_4+e_3+e_2+e_1.
\end{align*}
The negative lower-level coefficients do not affect domination: the
coefficient of the highest basis element $e_{n+1}$ is $1$.

For additional coefficient information, the $k$th coefficient of $f$ is
\[
 [Z^k]f=(-1)^k\sum_{1\le j_1<\cdots<j_k}
                         t^{e_{j_1}+\cdots+e_{j_k}},
 \qquad
 v\bigl([Z^k]f\bigr)=e_1+\cdots+e_k,
\]
the minimum coming from the unique subset $\{1,\ldots,k\}$, whose coefficient
is $(-1)^k$ and therefore cannot cancel. Division by the finite integer $k$
still leaves a cofinal sequence of Archimedean classes, which independently
exhibits the coefficient growth required by \cref{ent:thm:criterion}. The
reciprocal $1/\sigma_j$ has the explicit support-certified form
\[
 \frac1{\sigma_j}
   =t^{e_j}\sum_{r\ge0}(-1)^r t^{r(e_j+e_{j+1})},
\]
a valid positive-support geometric identity by
\cref{ent:lem:neumann}(c). Its support need not itself be cofinal in the
rank-infinite group; that is irrelevant to the legitimacy of this single
coefficient. Entirety of the product is supplied by the cofinal escape of the
factor indices $j$, not by a false claim about $r(e_j+e_{j+1})$.

\subsubsection*{The same workspace seen from its canonical product}
The natural product over the \emph{whole} basis,
\begin{equation}\label{ent:eq:infinite-product}
 P(Z)=\prod_{j\ge0}(1-t^{e_j}Z),
\end{equation}
differs from \eqref{ent:eq:f-explicit} by the single linear factor of index
zero:
\begin{equation}\label{ent:eq:P-vs-f}
 P(Z)=(1-t^{e_0}Z)f(Z)=(1-\omega^{-1}Z)f(Z).
\end{equation}
It is therefore an equally explicit member of the same ring, with one extra
zero at $\omega$.

\begin{theorem}[A genuine infinite-rank surcomplex entire function]
\label{ent:thm:infinite-example}
The product \eqref{ent:eq:infinite-product} belongs to
$\mathcal E_{\Gamma_\infty}$ and is nonpolynomial. Its zeros are exactly
\begin{equation}\label{ent:eq:surreal-zeros}
 r_j=t^{-e_j}=\omega^{\omega^j}\qquad(j\ge0),
\end{equation}
all simple and positive surreal. Its coefficient of $Z^k$ is
\begin{equation}\label{ent:eq:elementary-coeff}
 c_k=(-1)^k\sum_{j_1<\cdots<j_k}t^{e_{j_1}+\cdots+e_{j_k}},
 \qquad
 v(c_k)=e_0+e_1+\cdots+e_{k-1}
\end{equation}
for $k\ge1$, with $c_0=1$. At a scale $s\in\Gamma_\infty$ put
\[
 N_s=\#\{j:e_j+s<0\},\qquad M_s=\#\{j:e_j+s=0\}.
\]
These are finite, $M_s\in\{0,1\}$, and
\begin{align}
 m_s(P)&=\sum_{e_j+s<0}(e_j+s),\label{ent:eq:product-profile}\\
 I_s(P)(X)&=(-1)^{N_s}X^{N_s}(1-X)^{M_s}.\label{ent:eq:product-initial}
\end{align}
Finally
\begin{equation}\label{ent:eq:product-derivative}
 v\bigl(P'(r_j)\bigr)=e_j+\sum_{k<j}(e_k-e_j).
\end{equation}
\end{theorem}
\begin{proof}
Since $e_j\to+\infty$ cofinally, the proposed zero divisor is radially
finite, and \cref{ent:thm:product} gives the product's existence, entireness
and exact simple zero set. Expanding over finite subsets gives
\eqref{ent:eq:elementary-coeff}; among the $k$-element subsets, the sum of
the first $k$ basis elements is uniquely least in the largest-index order and
its coefficient is $(-1)^k$, so it cannot cancel and every $c_k$ is nonzero.

The slope sequence $(e_0+\cdots+e_{k-1})/k$ is cofinal: once $k-1$ exceeds
the largest index occurring in a prescribed group element, its positive
$e_{k-1}$ coordinate dominates that element. This is a second, direct
verification of entireness via \cref{ent:thm:criterion}.

At scale $s$ only finitely many factors have $e_j+s\le0$. A factor of
negative value contributes that value to $m_s$ and contributes $-X$ to the
normalized initial polynomial; a factor of value zero contributes $1-X$; all
positive factors contribute $1$. This proves
\eqref{ent:eq:product-profile}--\eqref{ent:eq:product-initial}.

Differentiating the factor of index $j$ and evaluating at $r_j$ kills all
other product-rule terms, giving
\[
 P'(r_j)=-t^{e_j}\prod_{k\ne j}(1-t^{e_k-e_j}).
\]
The factors with $k>j$ form a positive-support unit, and the finitely many
factors with $k<j$ have valuations $e_k-e_j$; adding these proves
\eqref{ent:eq:product-derivative}.
\end{proof}

This example is confined to no single Archimedean scale, yet its entire
theory is nontrivial and its zero divisor completely explicit. Together with
\cref{ent:rem:two-existence} it shows that nonpolynomial entire functions and
nonconstant restricted units are genuinely different existence questions.

\subsection{Two series with the same leading valuations and different
boundary behaviour}
\label{ent:subsec:hidden-tails}
Let $\Gamma=\Q\epsilon\oplus\Q\eta$ with $0<\epsilon\ll\eta$ and define
\begin{align}
 A(Z)&=\sum_{n\ge1}t^{n^2\epsilon}Z^n,\label{ent:eq:hidden-A}\\
 B(Z)&=\sum_{n\ge1}
       \left(t^{n^2\epsilon}+t^{\eta-n^2\epsilon}\right)Z^n.
       \label{ent:eq:hidden-B}
\end{align}
Their nonzero coefficients have exactly the same valuations $n^2\epsilon$.
Nevertheless their \emph{monomial} strong evaluation domains differ. For
$s=a\eta+b\epsilon$ with $a,b\in\Q$,
\begin{equation}\label{ent:eq:hidden-profile}
 \begin{array}{c|ccc}
 &a<0&a=0&a>0\\\hline
 A(t^s)&\text{not strong}&\text{strong}&\text{strong}\\
 B(t^s)&\text{not strong}&\text{not strong}&\text{strong}
 \end{array}
\end{equation}
For $a>0$ the leading valuations have $\eta$ coordinate $na\to+\infty$, and
\cref{ent:lem:cofinal-tails} proves both positive assertions. For $a<0$ those
leading valuations strictly decrease. For $a=0$ the exponents of $A(t^s)$ are
$(n^2+nb)\epsilon$, which after finitely many terms strictly increase with
only finite repetition at any exponent, so they form a well-ordered support.
The extra exponents of $B(t^s)$ are $\eta+(nb-n^2)\epsilon$, which strictly
decrease for all sufficiently large $n$; they destroy strong summability
without affecting a single leading coefficient valuation.

Neither series is entire over the whole rank-two field: their coefficient
slopes $n\epsilon$ remain below $\eta$, so \cref{ent:thm:criterion} correctly
excludes both. There is no tension between this example and the
valuation-only criterion. Leading valuations can fail to decide strong
summability at a \emph{boundary} scale, while global summability at
\emph{every} scale removes the ambiguity. Formula
\eqref{ent:eq:hidden-profile} is stated for monomial arguments only; it does
not silently assert a classification for arbitrary nonmonomial arguments of
$B$.

\subsection{One extension preserves the function; another destroys entireness}
\label{ent:subsec:two-extensions}
Adjoin $\omega^{1/2}$ to $\Gamma_\infty$ and take the divisible additive span.
This is an ordered-group extension in which $\Gamma_\infty$ remains cofinal,
since a sufficiently large integer power of $\omega$ dominates every finite
combination of the added and old generators. The series $P$, and likewise the
pair $f,g$, remain entire by
\cref{ent:cor:entire-extension,ent:cor:no-repair}; in particular the B\'ezout
failure persists.

Instead adjoin the new generator $e_*=\omega^\omega>\Gamma_\infty$ and set
$\Delta=\Gamma_\infty\oplus\Q e_*$. At the point $Z=t^{-e_*}$ the nonzero
summands of $P$ have leading exponents
$(e_0+\cdots+e_{k-1})-ke_*$, whose consecutive differences are
$e_k-e_*<0$; the defining family is therefore not strongly summable, exactly
as \cref{ent:thm:extension-criterion} predicts.
\Cref{ent:thm:no-continuation} rules out any alternative entire power series
over $K_\Delta$ agreeing with $P$ on the old field, or even on $\C$. On the
still-defined domain $\mathcal D_{\Gamma_\infty,\Delta}$, however, $P$ is
unchanged and its zeros remain exactly \eqref{ent:eq:surreal-zeros}, by
\cref{ent:thm:extension,ent:thm:zero-conservation}. This concrete pair is the
clearest illustration of \cref{ent:rem:reconcile}.

\subsection{Why finite rank does not give the counterexample}
\label{ent:subsec:finite-rank}
For $\Gamma=\Q^d$ in lexicographic order let $h$ lie in the largest
Archimedean class. Integer multiples of $h$ are cofinal, so $h$ is an order
unit and the entire ring is B\'ezout by \cref{ent:thm:bezout}, even when
$d>1$. The same holds for $\Q e_1\oplus\Q e_2$ with $e_2\ggA e_1$: take
$h=e_2$. It is \emph{absence of a greatest Archimedean class}, together with
countable cofinality, that forces the negative case --- not high rank.

\subsection{An uncountable hierarchy gives polynomial rigidity}
\label{ent:subsec:omega-one}
Let
\begin{equation}\label{ent:eq:Gamma-omegaone}
 \Gamma_{\omega_1}=\bigoplus_{\alpha<\omega_1}\Q e_\alpha,
 \qquad e_\alpha=\omega^\alpha,
\end{equation}
again ordered by the largest nonzero index. Its cofinality is $\aleph_1$: the
basis is cofinal, whereas any countable family of finite supports has its
indices bounded below $\omega_1$ and a larger basis element bounds the whole
family. Therefore
\begin{equation}\label{ent:eq:uncountable-rigidity}
 \mathcal E_{\Gamma_{\omega_1}}
 =\mathcal T_{\Gamma_{\omega_1}}
 =K_{\Gamma_{\omega_1}}[Z],
\end{equation}
the middle equality because a nonpolynomial restricted series would itself
supply a countable sequence of coefficient valuations tending cofinally to
infinity. This example shows that infinite rank by itself does not even
guarantee nonpolynomial entire functions: rank is not the driver, the order
type above the scales is.

\subsection{The separation table}
\label{ent:subsec:table}
\begin{center}\small
\begin{tabularx}{\textwidth}{@{}p{0.24\textwidth}cYYY@{}}
\toprule
Value group & $\cf$ & Has an order unit & Nonpolynomial entire series &
Nonconstant restricted units\\\midrule
$\Q$ or $\R$ & $\aleph_0$ & Every positive element & Yes & Yes\\[0.45em]
$\Q\epsilon\oplus\Q\eta$, $\epsilon\ll\eta$ & $\aleph_0$ &
Yes; not every positive element & Yes & Yes, with the sharper gap
test\\[0.45em]
$\Gamma_\infty$ & $\aleph_0$ & No & Yes & No\\[0.45em]
$\Gamma_{\omega_1}$ & $\aleph_1$ & No & No & No\\
\bottomrule
\end{tabularx}
\end{center}

The second and third rows are the two nonpolynomial cases of
\cref{ent:thm:main}; the third is the non-B\'ezout one. Note that $\R$ is
uncountable of rank one with $\cf=\aleph_0$, while $\Gamma_\infty$ has
infinite rank with $\cf=\aleph_0$: neither cardinality nor rank is the
invariant at work.

\section{Relation to other reports in this collection}
\label{ent:sec:collection}

This report is a merge of two independently written manuscripts. Both were
prepared against pinned snapshots of the same repository, and both state
explicitly which of their neighbours' results they do \emph{not} claim. Those
statements are preserved here, and the cross-references are reciprocal.

\paragraph{Rank-one Berkovich geometry --- this report extends it.}
The report \emph{From Surcomplex Hahn Fields to Berkovich Geometry}
\cite{RepoRankOne} fixes $K=\C((t^\R))$: the value group is $\R$, the field
is spherically complete, and $|a|_v=\exp(-v(a))$ is a genuine real-valued
absolute value. It is, in its own words, the one place in the collection
where convergence is real convergence. Its entire-function section proves
zero-free entire rigidity and a convergent canonical product \emph{there},
and neither source manuscript merged here claims to originate either point.
The present report is the arbitrary-rank generalization. Three specific
relations should be recorded.
\begin{itemize}[leftmargin=1.4em]
\item \emph{Where rank one sits in the classification.} For $\Gamma=\R$ we
have $\cf(\R)=\aleph_0$ and $1$ is an order unit, so rank one falls squarely
in case (ii) of \cref{ent:thm:main}: nonpolynomial entire functions exist,
$\mathcal E_\R$ is a non-Noetherian B\'ezout domain, and arbitrary radially
discrete Hermite interpolation is available. The pathology of case (iii) is
invisible at rank one.
\item \emph{Why rank one is not general.} That report's
\texttt{prop:rankobstruction} proves that $\Q+\Q\omega$ admits no
order-preserving \emph{injective} additive map to $\R$. That is exactly why
\cref{ent:lem:coarsening} constructs a \emph{coarsening} with a possibly
nonzero kernel $H$ rather than an embedding, and why it is available only
when an order unit exists. The two statements are complementary: one forbids
faithful scalarization in general, the other supplies the best available
non-faithful substitute in the order-unit case.
\item \emph{The ring chains resemble each other and are not equal.} That
report's strict chain $K[Z]\subsetneq\mathcal T\subsetneq\mathcal P
\subsetneq\mathcal A\subsetneq\mathcal F\subsetneq\mathcal E^-$ locates its
polynomial-coefficient and formal algebras. Our chain
\eqref{ent:eq:chain} resembles it but is not the same chain, and the symbols
$\mathcal T$, $\mathcal A$, $\mathcal E$ do not denote the same rings in the
two reports. \Cref{ent:conv:ring} forbids transferring a statement across
that boundary.
\end{itemize}
Conversely, the rank-one report's own entire-function section should be read
with the knowledge that the arbitrary-rank classification now exists and that
its rank-one instance lands in the B\'ezout case.

\paragraph{Analysis and Foundations --- this report is published beside them.}
\Cref{ent:cor:all-no} re-derives a conclusion the collection already contains
twice: the \emph{Analysis} report proves all-scale polynomial rigidity
directly (its Theorem \texttt{f:thm-rigidity}), and the \emph{Foundations}
report records the exact boundary of the phenomenon with a pair of series
over $\C((t^\Q))$ that differ only in the sign of an exponent, noting that
$\sum t^{+n^2}z^n$ is strongly summable at every point of that field and is
\emph{nonpolynomial} while remaining coherent at every radius, and that this
does not refute the all-scale theorem but shows that the all-scale theorem
does not transfer to a \emph{fixed} value group. That example is precisely
\cref{ent:rem:criterion-strength} of this report, and the trichotomy of
\cref{ent:thm:main} explains why it works over $\Q$ --- $\cf(\Q)=\aleph_0$
and $1$ is an order unit, so $\Q$ is a case (ii) group --- and what it would
take to break it, namely replacing $\Q$ by a group of uncountable cofinality.
Conversely, our reduction in \cref{ent:cor:all-no} consumes the foundations
report's \emph{localization principle}: every set of surcomplex numbers lies
in one divisible set-sized workspace. That principle is what makes the
reduction legitimate, and it is cited here, not reproved.

\paragraph{Global divisors and analytic geometry --- different rings.}
The global-divisor report \cite{RepoDivisors} works in
$\mathcal O(U)((t^\Gamma))$, with ordinary holomorphic coefficient functions
on a common complex domain; its divisor, interpolation and Picard
obstructions live in a different function ring from the strongly entire ring
studied here, and nothing in this report contradicts or restates them. The
analytic-geometry report states the ring discipline that
\cref{ent:conv:ring} invokes by name, for its own four coefficient rings.
Neither the objects nor the theorems transfer.

\section{What the results do and do not establish}
\label{ent:sec:scope}

\subsection{Three independent distinctions}
\label{ent:subsec:three-distinctions}
The first distinction is between strong Hahn summation and intrinsic
valuation convergence. Positive well-ordered support suffices for the former;
cofinal progress through every group threshold is needed for the latter. On a
single disk they differ. For a series strongly evaluable at \emph{every}
point, \cref{ent:thm:criterion} shows the distinction disappears within its
original field.

The second is between countable cofinality and an order unit. Countable
cofinality supplies a sequence of ever larger scales, possibly moving through
infinitely many Archimedean classes; an order unit supplies all those bounds
by the integer multiples of one element. Entire functions need only the
first; a nonconstant restricted inverse requires the second. The group
$\Gamma_\infty$ separates them, and that separation is the whole content of
the gap between cases (ii) and (iii) of \cref{ent:thm:main}.

The third is between an intrinsic topology in one workspace and the fine
topology tested against every surreal scale. The same countable sequence may
converge intrinsically in $\K$ and fail to converge after a noncofinal
enlargement. None of our proofs identifies these topologies. See
\cref{ent:subsec:three-meanings}.

\subsection{Established, and no more}
The results isolate three logically different phenomena. Coefficient slopes
determine whether an ordinary countable power series can be evaluated at
every scale in a fixed field. Support-controlled preparation and canonical
products completely describe zeros and divisibility within that category. The
order structure above those scales determines whether arbitrary values and
jets can be prescribed and whether coprime functions generate the unit ideal.

In particular the factorization monoid alone cannot decide the B\'ezout
property. In both nonpolynomial cases a function up to a scalar is exactly a
radially finite zero divisor, and every effective subdivisor is again
principal; nevertheless only one of those cases admits unrestricted discrete
interpolation. The additive structure of the ring, not merely its
multiplicative divisibility relation, is essential.

The non-B\'ezout example is not a delicate choice of a transcendental
ordinary complex function or of an uncomputable coefficient. Its nodes,
displacements and factor coefficients are explicit Hahn expressions over
rational data. What cannot be replaced by a finite calculation is the
quantification over all coefficient indices and all Archimedean scales.

\subsection{Scope that must not be silently enlarged}
\label{ent:subsec:scope-limits}
The following restrictions are stated by the source manuscripts and are
retained without weakening.

\paragraph{Hypotheses that cannot be relaxed.}
All supports and index sets are sets; there is no proper-class summation.
The field is a \emph{full} Hahn field with complex residue coefficients, and
$\Gamma$ is divisible. Algebraic closedness is used essentially, to split the
prepared finite polynomial; it is the sole import of
\cref{ent:thm:zero-conservation}. A nondivisible group, or a
support-restricted Hahn subfield, may lose the roots needed to split the
prepared polynomial, so the factorization theorem cannot simply be copied
there. The Hahn--Neumann support lemmas
(\cref{ent:lem:neumann}) and algebraic closedness \cite{Poonen} are imported
inputs, not proved here.

\paragraph{Categories excluded.}
The ring is one-variable and strong-summation entire. It is \emph{not} the
ring of ordinary-holomorphic-coefficient Hahn families on one common complex
domain, not a ring of arbitrary fine-local surcomplex germs, and not a
several-variable analytic algebra. The arbitrary-rank preparation theorem is
proved in $\C[X]((t^\Gamma))$, not asserted in a larger power-series
coefficient ring. There is no contour integral here, no global surcomplex
exponential at unrestricted imaginary arguments, and nothing about the
Berarducci--Mantova derivation; the statement about primitives in
\cref{ent:prop:operations} concerns termwise $Z$-integration with
coefficients fixed and says nothing about differential equations in the
coefficient field. All entire functions here are represented by one ordinary
formal power series at the origin: translation preserves this class, but does
not identify it with every locally analytic or surreal-differentiable
function. Piecewise functions on fine-topological domains, Gonshor's
exponential, locally represented classes and other support or summation
conventions lie outside the hypotheses, so
\cref{ent:cor:all-no} is \emph{not} a general no-go theorem for surreal
analysis.

\paragraph{Limits of the coarsening and of the extension theorem.}
The rank-one coarsening of \cref{ent:lem:coarsening} changes residue
information even though it preserves the topology and the entire ring;
residue-direction data is not identified across the two valuations. The
scalar-extension results keep the coefficient field $\C$ fixed and use
canonical ordered-group inclusions of full Hahn fields only: they are not
theorems about arbitrary extensions of incomplete subfields, about changing
residue fields, or about analytic sheaves. The hidden-tail table
\eqref{ent:eq:hidden-profile} is stated for monomial arguments and does not
classify arbitrary nonmonomial arguments of $B$. The single-variable
uniform-domain statement is \emph{false} in several variables --- a series
independent of one coordinate imposes no constraint in that coordinate ---
and the correct multivariable replacement is posed below as an open question.
No density of $\C$ in any fine surreal topology is asserted, and strong Hahn
summability, intrinsic valuation convergence and the fine surreal topology
are deliberately never identified.

\paragraph{What is and is not claimed about priority and verification.}
The results of this report are self-declared proposed original contributions,
not a solution of any named published conjecture. Both source manuscripts
state that priority is \emph{not} certified and that the literature checks
were targeted rather than exhaustive; in one of them several general web
searches returned irrelevant results and were explicitly not treated as
evidence of absence, and in the other a connector keyword search returning
nothing was likewise not treated as proof. The rank-one factorization
specialization is stated to be classical; standard Hahn algebra, algebraic
closedness and support lemmas are not claimed; Cherry himself calls his
factorization results well known; the rank-two distinction between positive
valuation and topological nilpotence is established
\cite[Section~6.2]{Conrad}; and the all-surcomplex polynomial corollary is
credited to the repository, as recorded in \cref{ent:rem:rigidity-credit}.
One source's audit further notes that Lazard was checked bibliographically
only and that Conway and Gonshor were not re-audited in full; its repository
review was a targeted read of four files, not a line-by-line proof audit or
an independent build. Neither manuscript was independently refereed and
neither was machine-checked: one records \texttt{formal\_verification false},
\texttt{independent\_peer\_review false} and \texttt{priority\_certified
false} in its build record, and the other states flatly that no Lean
verification is claimed. No theorem below depends for its validity on the
correctness of an unrefereed repository report. The repository's Lean modules
on finite algebra and Hahn supports are plausible \emph{prerequisites} for a
future formalization, not a certificate that these proofs have been
formalized; missing implementation is not treated here as an open
mathematical problem. No repository files under \texttt{sources/},
\texttt{code/} or \texttt{data/} have been modified.

\paragraph{What the finite checks do not do.}
See \cref{ent:app:checks} for the full statement. In summary: the two
delivered suites verify $831$ and $664$ finite exact assertions respectively,
and they prove \emph{no} infinite theorem. Both articles and both READMEs
state that the finite truncations of the counterexample pair \emph{are}
coprime polynomials and \emph{do} satisfy polynomial B\'ezout identities, so
no finite computation can witness \cref{ent:thm:no-bezout}. The preparation
checks run modulo $t^9$ in an integer exponent subgroup, not an
arbitrary-rank group, and the finite dominance samples test only
$k\in\{1,2,10,100,10000\}$, not all integer multiples. The checks cannot
verify arbitrary well-ordered supports, the cofinality dichotomy, the
necessity direction of \cref{ent:thm:units}, or any infinite product
identity.

\subsection{Questions opened by the classification}
\label{ent:subsec:questions}
The following are research questions generated by this work, not claims that
any particular publication lists them as open.

\begin{question}[Exact interpolation image]\label{ent:q:image}
When $\Gamma$ has countable cofinality but no order unit, characterize the
image of
\[
 \E\longrightarrow\prod_{x\in\supp D}\K[U]/(U^{D(x)})
\]
for a prescribed radially finite divisor $D$.
\end{question}
This question is open on the article's own terms. The growth barrier gives
explicit necessary inequalities and proves that the image can be proper, but
sufficiency for arbitrary nodes and jets is not established. A satisfactory
answer should encode the compatibility of divided differences across convex
subgroups, rather than only the size of individual values.

\begin{question}[Finitely generated ideals beyond their common divisor]
\label{ent:q:ideals}
Can one describe the ideal $(f,g)$ by its common zero divisor together with
an explicit obstruction in an interpolation quotient?
\end{question}
There is already an exact reduction: after dividing out the GCD, the B\'ezout
problem is whether the inverse jets of one function along the other's zero
divisor lie in the interpolation image. What is missing is a usable intrinsic
description of that image. \Cref{ent:thm:no-bezout} proves that a
zero-divisor description alone is insufficient.

\begin{question}[Several variables]\label{ent:q:several}
Which parts of the cofinality criterion and of the B\'ezout obstruction
extend to several ordinary variables, and which local preparation statements
retain explicit Hahn support certificates? In particular, since the
single-variable uniform extension domain is false in several variables, what
is the correct replacement --- presumably one depending on the unbounded
degree directions of the support, including the boundary where higher Hahn
tails matter?
\end{question}
One must separate the cofinal growth of total degree from the geometry of
positive-dimensional zero sets. Cherry's several-variable work is an
important rank-one antecedent \cite{Cherry}; no arbitrary-rank multivariable
GCD theorem is established here.
\Cref{ent:app:multivariate} proves only the finite-variable convergence and
unit criteria, and it does not extend the one-variable divisor
classification to higher-dimensional zero sets.

\begin{question}[Restricted and incomplete subfields]\label{ent:q:subfields}
How much of the divisor theory survives for incomplete surreal subfields or
for support-restricted Hahn subfields? Our proofs isolate two requirements:
closure under the positive-support preparation construction, and enough
algebraic roots to split the finite prepared polynomials. A classification of
natural surreal subfields meeting both would turn these full-Hahn results
into a more computationally selective theory.
\end{question}

\subsection{A proof-assistant route}
The research statement does not require a formalization of the entire proper
class $\No$. A first formal development can work over a set-sized divisible
linearly ordered abelian group and a Hahn field. The most useful first target
is the global criterion \cref{ent:thm:criterion}, whose nontrivial
implication is a short ordered-group argument plus the definition of strong
summability; the next is the coarsened polynomial argument for
\cref{ent:thm:units}. Support preparation should be formalized only after the
finite-word Neumann lemma and arbitrary regrouping are available. The
explicit obstruction then requires only ordered-group inequalities, finite
product valuations and evaluation homomorphisms. This is a dependency plan.
It is not a claim that any of the proofs has been checked in Lean, and no
Lean code or kernel-check claim is included in this package.

\section{Conclusion}
\label{ent:sec:conclusion}
The threshold for nonpolynomial entire behaviour in a Hahn--surcomplex
workspace is countable cofinality, not rank one. The threshold for a
B\'ezout identity between coprime entire functions, and equally for a
nonconstant restricted inverse, is strictly stronger: a positive coefficient
gap must generate a cofinal sequence of multiples. These two conditions
separate in an explicit infinite-rank subgroup of the surreal numbers, and
the separation is witnessed by two entire functions with disjoint simple zero
sets whose ideal is proper.

After enlarging the value group, every nonpolynomial entire series has the
same exact strong domain: the valuation ring of the convex hull of its
original scales. Cofinality is precisely what preserves entireness, while
support preparation preserves every existing zero and prevents new ones
throughout the surviving domain. The canonical product theorem then shows
that coefficients, zeros and the extension boundary encode the same cofinal
scale structure. Together this gives one explanation of both nonpolynomial
fixed-workspace functions and polynomial rigidity on the whole surcomplex
class, without confusing Hahn identities with convergence in the fine surreal
topology.

\appendix
\section{Proof-dependency and hypothesis audit}
\label{ent:app:audit}

\subsection{The dependency chain}
The main proof can be read along the following chain.
\[
 \begin{gathered}
 \text{Hahn support calculus}\\
 \Downarrow\\[-2pt]
 \text{cofinal coefficient criterion and the algebra }\A\\
 \Downarrow\\[-2pt]
 \text{support preparation and finite root counts}\\
 \Downarrow\\[-2pt]
 \text{canonical products, divisibility, and GCDs}\\
 \Downarrow\\[-2pt]
 \begin{array}{c@{\qquad\qquad}c}
 \text{new Archimedean scales}&\text{one order unit}\\
 \Downarrow&\Downarrow\\
 \text{growth barrier and no B\'ezout}
 &\text{Hermite interpolation and B\'ezout}
 \end{array}
 \end{gathered}
\]
The arrow from preparation to root counting additionally uses algebraic
closedness of the Hahn field. The arrow to interpolation uses completeness of
the real-valued coarsening, proved directly in \cref{ent:lem:coarsening}. No
proof of the negative result depends on the positive interpolation theorem.

The scalar-extension results of \cref{ent:sec:extension} sit on a deliberately
separate branch. The restricted-unit theorem
(\cref{ent:thm:units}) uses Gauss multiplicativity
\eqref{ent:eq:gauss-mult} and a convex coarsening only; it uses neither zero
counting nor canonical products. The universal-domain theorem
(\cref{ent:thm:extension}) uses the global criterion and integral evaluation;
it does not use factorization. \emph{Only} zero conservation
(\cref{ent:thm:zero-conservation}) imports both algebraic closedness and
preparation. The divisor classification is downstream of the local tools and
the elementary entire algebra. Thus the principal classifications are not
consequences of an unproved global factorization assertion, and conversely
the factorization proof never assumes that an arbitrary zero-free restricted
series is restricted-invertible --- by
\cref{ent:ex:ranktwo-unit} that is exactly the assertion that can fail in
higher rank.

\subsection{Where an attained minimum is load-bearing}
\paragraph{Attained minima.}
The proof of \cref{ent:thm:criterion} uses a minimum attained at an ordinary
\emph{finite} coefficient index, not a termwise inequality. The proof of
\cref{ent:lem:growth-barrier} uses an attained minimum in the same way:
termwise strict inequalities alone would not justify a strict inequality for
an arbitrary limiting process. These are the two places in the article where
the distinction is load-bearing, and neither can be replaced by an estimate
valid only term by term.

\paragraph{Two independent infinities.}
The finite index $k$ in a power series is not a Hahn exponent. Integer
multiples of a positive exponent need not be cofinal. The support lemma
\cref{ent:lem:neumann}(b) handles geometric-series expansions without that
property; the positive interpolation proof uses it only after constructing a
legitimate real-valued coarsening, which is available only under the
order-unit hypothesis.

\paragraph{Divisibility is not assumed in its own proof.}
The quotient $F/G$ in \cref{ent:thm:factorization} is first formed, after
removing the shared zero at the origin, in $\K[[Z]]$. Only preparation at
every scale proves it entire. This avoids assuming the global divisibility
theorem in its own proof.

\paragraph{The counterexample is irreducibly infinite.}
No finite truncation of the pair \eqref{ent:eq:paired-products} can witness
non-B\'ezout behaviour: the finite truncations are polynomials with disjoint
root sets and therefore generate the unit ideal. The contradiction requires a
single candidate coefficient lower bound to fail at cofinally many
successively dominating scales.

\paragraph{Why preparation must be a support recursion.}
\Cref{ent:cor:prep-extension} --- the fact that the prepared factorization is
unchanged by enlarging the value group --- is available precisely because
preparation was proved by well-founded recursion over $S^*$. An iteration
whose convergence had been asserted in $\K$ would have to have that
convergence re-established in $\KD$, and \cref{ent:ex:ranktwo-unit} shows it
can genuinely fail.

\section{Finite-variable convergence and unit criteria}
\label{ent:app:multivariate}
Let $d\ge1$ be finite, write $\alpha=(\alpha_1,\ldots,\alpha_d)\in\N^d$ and
$|\alpha|=\alpha_1+\cdots+\alpha_d$. Define
\[
 \mathcal T_{\Gamma,d}
 =\left\{\sum_\alpha a_\alpha X^\alpha:
              v(a_\alpha)\to+\infty\text{ as }|\alpha|\to\infty\right\},
\]
and define $\mathcal E_{\Gamma,d}$ by strong summability of the \emph{whole}
family $(a_\alpha x^\alpha)_{\alpha\in\N^d}$ for every $x\in\K^d$. Summing
first along one coordinate and then another is not the definition.

\begin{theorem}\label{ent:thm:multi-criterion}
For a finite-variable formal series $F=\sum_\alpha a_\alpha X^\alpha$,
\[
 F\in\mathcal E_{\Gamma,d}
 \quad\Longleftrightarrow\quad
 \frac{v(a_\alpha)}{|\alpha|}\to+\infty
 \quad(|\alpha|\to\infty,\ |\alpha|>0).
\]
It suffices to test strong summability at the diagonal monomial points
$(t^s,\ldots,t^s)$ for all $s\in\Gamma$. Nonpolynomial members exist exactly
when $\cf(\Gamma)=\aleph_0$.
\end{theorem}
\begin{proof}
At a diagonal monomial point the summand values are
$v(a_\alpha)+|\alpha|s$, and there are only finitely many multi-indices of
any bounded total degree. If these values do not tend cofinally to infinity,
strong summability gives a common lower bound and an infinite bounded-above
subfamily with strictly increasing total degrees. Subtracting
$|\alpha|\delta$, for a positive $\delta$ larger than the length of the
bounding interval, produces the same strictly descending leading-exponent
sequence as in \cref{ent:thm:criterion}. Thus diagonal strong summability
implies the slope criterion.

Conversely, for a point with all coordinates nonzero set
$s=\min_i v(x_i)$; then
$v(a_\alpha x^\alpha)\ge v(a_\alpha)+|\alpha|s\to+\infty$, and because
bounded total degrees contain finitely many indices,
\cref{ent:lem:cofinal-tails} applies after any enumeration respecting total
degree. If some coordinates vanish, discard the vanishing terms and apply the
same argument to the remaining coordinates. The cofinality assertion follows
exactly as in one variable; for existence use an entire nonpolynomial in
$X_1$ alone.
\end{proof}

\begin{theorem}\label{ent:thm:multi-unit}
A nonconstant $F\in\mathcal T_{\Gamma,d}$ is a unit of that algebra if and
only if $a_0\ne0$ and
$\min_{\alpha\ne0,\ a_\alpha\ne0}v(a_\alpha/a_0)$ is a positive order unit of
$\Gamma$.
\end{theorem}
\begin{proof}
Regrouping into $\C[X_1,\ldots,X_d]((t^\Gamma))$ gives multiplicativity of
the Gauss valuation. If $FG=1$, normalize both constant coefficients to $1$;
their Gauss values are nonpositive and sum to zero, so their initial
polynomials multiply to $1$ and are constant, and all nonconstant coefficient
values are positive. If their minimum $\delta$ is not an order unit, coarsen
by the proper convex subgroup \eqref{ent:eq:Hdelta} it generates:
restrictedness reduces both series to finite polynomials over the coarsened
residue field, and at least one nonconstant coefficient of $F$ survives,
contradicting invertibility of its reduction. Conversely an order-unit gap
makes the geometric inverse restricted, with precisely the tail estimates of
\cref{ent:thm:units}; finite-variable convolution has finitely many terms at
each multi-index, so those estimates are unchanged.
\end{proof}

Finite-dimensionality is used in the finiteness of bounded-degree index sets.
These statements assert nothing for infinitely many variables, and they do
not extend the one-variable divisor classification to higher-dimensional zero
sets. As noted in \cref{ent:subsec:scope-limits}, the single-variable uniform
extension domain of \cref{ent:thm:extension} is \emph{false} in several
variables.

\section{Reproducible finite checks}
\label{ent:app:checks}
Two independent exact-arithmetic suites accompany this report, preserved
unmodified in \texttt{code/} and \texttt{data/}. Both use only the Python
standard library, with exact \texttt{fractions.Fraction} arithmetic and no
floating point, and both represent group elements as tuples of rational
coordinates ordered by the largest nonzero coordinate.

\paragraph{Suite A (831 checks).}
\texttt{verify\_examples.py} checks finite instances of the
reverse-lexicographic exponent inequalities, the valuation formula
\eqref{ent:eq:lambda} in the first nine coordinates of $\Gamma_\infty$,
leading exponents of finite canonical products, a finite-truncation
implementation of the preparation recursion
\eqref{ent:eq:prep-rec1}--\eqref{ent:eq:prep-rec2}, and the local inverse-jet
identity underlying \eqref{ent:eq:interpolation-r}. The delivered run passed
all \textbf{831} checks, with category counts $72$ paired-root checks, $40$
finite dominance samples, $88$ product-coefficient checks, $456$ preparation
checks and $175$ inverse-jet checks; the record is the delivered
verification JSON. The script takes a \emph{required} output path argument
and exits nonzero without one; this is a usage requirement designed so that a
re-run cannot overwrite the delivered record, not a defect. Run correctly it
reproduces $831$ of $831$.

\paragraph{Suite B (664 checks).}
\texttt{verify.py} checks finite truncations of
\eqref{ent:eq:infinite-product} for $N=1,\ldots,8$ --- each coefficient's
leading exponent, its binomial subset count and its signs --- exactness and
simplicity of each root $t^{-e_j}$ and the simple-root derivative
\eqref{ent:eq:product-derivative}, closed- and open-ball Newton counts at
sampled scales including half-shifts by $e_0$, descent of the valuations
after adjoining the dominant scale $e_9$, the rank-two hidden-tail comparison
of \cref{ent:subsec:hidden-tails} for $n=1,\ldots,100$ and
$\lambda\in\{-7,-\tfrac12,0,\tfrac32,10\}$, and the preparation recursion for
\eqref{ent:eq:cubic} through perturbation order eight. The delivered run
reports \textbf{664} passing checks.

\paragraph{What the checks cannot do.}
Both suites print their own scope disclaimers, and both are correct to do so.
The counts are counts of finite assertions, not confidence scores for any
infinite theorem. Neither suite can prove cofinality statements, Hahn support
well-ordering for every permitted group, any infinite product identity, the
necessity direction of \cref{ent:thm:units}, or the nonexistence of a global
B\'ezout solution. In the preparation checks the $t$-adic truncation is
explicitly finite, so higher coefficients are discarded only after the
claimed congruence has been specified, and the exponent subgroup used is an
integer one rather than an arbitrary-rank group. The finite dominance samples
test only $k\in\{1,2,10,100,10000\}$. Most importantly, the finite
truncations of \eqref{ent:eq:paired-products} \emph{are} coprime polynomials
and \emph{do} satisfy polynomial B\'ezout identities, so no finite
computation can witness \cref{ent:thm:no-bezout}. What the checks can detect
is an indexing error, a sign error in a paired root valuation, a reversed
order convention, a wrong root-count endpoint or a wrong finite preparation
coefficient. Everything else rests on the proofs in the body.

\section{Source and novelty ledger}
\label{ent:app:sources}
This report merges two manuscripts, both dated 21 September 2026, each of
which pinned the repository at its own snapshot: one at commit
\texttt{4896a2808ce30e01b1c86ae3ba2295a64762d246} and the other at
\texttt{aa846271b4dcae2c055b216126a87210292ec19b}. The first re-read four
files at its pinned commit --- the root README, the documentation catalogue,
and the READMEs of the rank-one and global-divisor reports --- after its
working branch changed during preparation; the second cited the documentation
catalogue and the rank-one article. Both reviews were targeted reads, not
line-by-line proof audits or independent builds. Neither reinterprets missing
Lean implementation as an unsolved mathematical problem.

The primary literature checked across the two manuscripts comprises Poonen's
Hahn-field construction and algebraic-closedness corollary, Neumann's ordered
division-ring paper for the support calculus, Cherry's complete real-valued
non-Archimedean setting together with his one-variable factorization
statement and his lectures, Conrad's higher-rank valuation examples, Lazard's
bibliographic record as the historical reference for one-variable zeros, van
der Hoeven's generalized-series operator framework, and the Conway--Gonshor
surreal normal-form background with cross-checks against
\cite{ADH,BournezGuilmant}. Lazard was checked bibliographically only, and
Conway and Gonshor were not re-audited in full. General web searches were
attempted; several returned irrelevant results and were not treated as
evidence that a theorem is absent from the literature.

The strongest proposed original statements are
\cref{ent:thm:no-bezout} and its use in \cref{ent:thm:main}, together with
\cref{ent:thm:extension,ent:thm:zero-conservation,ent:thm:no-continuation}
and their assembly in \cref{ent:thm:main-extension}. The sharp workspace
criterion and the persistence result are
\cref{ent:thm:extension-criterion,ent:cor:no-repair}; the restricted-unit
criterion \cref{ent:thm:units} is closely related to established topological
nilpotence distinctions, and it is its precise necessity-and-sufficiency
formulation for \eqref{ent:eq:Talg}, including groups with no order unit,
that is proposed here. The rank-one factorization formula is explicitly
classical. The all-surcomplex polynomial corollary is credited to the
pre-existing all-scale theme in the repository
(\cref{ent:rem:rigidity-credit}). A later discovery of an antecedent would
change the priority discussion, not the content or the hypotheses of the
displayed proofs.

\clearpage
\phantomsection\addcontentsline{toc}{section}{References}
\begingroup\raggedright
\begin{thebibliography}{99}
\bibitem{Conway}
J. H. Conway, \emph{On Numbers and Games}, London Mathematical Society
Monographs, vol.~6, Academic Press, 1976.

\bibitem{Gonshor}
H. Gonshor, \emph{An Introduction to the Theory of Surreal Numbers},
London Mathematical Society Lecture Note Series, vol.~110,
Cambridge University Press, 1986.

\bibitem{Neumann}
B. H. Neumann, On ordered division rings, \emph{Transactions of the American
Mathematical Society} \textbf{66} (1949), 202--252.
\href{https://doi.org/10.1090/S0002-9947-1949-0032593-5}{doi:10.1090/S0002-9947-1949-0032593-5}.

\bibitem{Poonen}
B. Poonen, Maximally complete fields, \emph{L'Enseignement
Math\'ematique} (2) \textbf{39} (1993), 87--106.
Section~3 and Corollary~4 are the principal inputs used here.
\url{https://math.mit.edu/~poonen/papers/amsval.pdf}.

\bibitem{Lazard}
M. Lazard, Les z\'eros d'une fonction analytique d'une variable sur un
corps valu\'e complet, \emph{Publications Math\'ematiques de l'IH\'ES}
\textbf{14} (1962), 47--75.
\href{https://doi.org/10.1007/BF02684326}{doi:10.1007/BF02684326}.
\url{https://www.numdam.org/item/PMIHES_1962__14__47_0/}.

\bibitem{Cherry}
W. Cherry, Existence of GCD's and factorization in rings of
non-Archimedean entire functions, \emph{Contemporary Mathematics}
\textbf{551} (2011), 57--69.
\url{https://arxiv.org/abs/1007.0984}.

\bibitem{CherryLectures}
W. Cherry, \emph{Lectures on Non-Archimedean Function Theory}, 2009,
arXiv:0909.4509. In particular the preparation and product results in
Lectures~2--3. \url{https://arxiv.org/abs/0909.4509}.

\bibitem{Conrad}
B. Conrad, \emph{Lecture 6: More constructions with Huber rings}, Stanford
Number Theory Learning Seminar on perfectoid spaces, 31 October 2014,
informal lecture notes. In particular Section~6.2 and Example~6.2.1.
\url{https://virtualmath1.stanford.edu/~conrad/Perfseminar/Notes/L6.pdf}.

\bibitem{vdH}
J. van der Hoeven, Operators on generalized power series,
\emph{Illinois Journal of Mathematics} \textbf{45} (2001), no.~4,
1161--1190.
\href{https://doi.org/10.1215/ijm/1258138061}{doi:10.1215/ijm/1258138061}.

\bibitem{ADH}
M. Aschenbrenner, L. van den Dries, and J. van der Hoeven,
On numbers, germs, and transseries, arXiv:1711.06936, revised version.
\url{https://arxiv.org/abs/1711.06936}.
Used for bibliographic cross-checks of the surreal background.

\bibitem{BournezGuilmant}
O. Bournez and Q. Guilmant, \emph{Surreal fields stable under exponential and
logarithmic functions}, 2022, arXiv:2201.08199.
\url{https://arxiv.org/abs/2201.08199}.

\bibitem{RepoRoot}
V. Reshetnikov, \emph{Surreal}: repository README and reported
formalization scope, inspected 21 September 2026.
\url{https://github.com/VladimirReshetnikov/Surreal}.
AI-assisted research collection; the reports are unrefereed drafts, not
independent peer-reviewed validation.

\bibitem{RepoCatalogue}
V. Reshetnikov, \emph{The surreal and surcomplex reports}, documentation
catalogue \texttt{docs/README.md} in the \emph{Surreal} repository,
inspected 21 September 2026.

\bibitem{RepoRankOne}
\emph{From Surcomplex Hahn Fields to Berkovich Geometry: rank-one
convergence, Tate algebras, analytic annuli, and an entire-function case
study}, research report in the \emph{Surreal} repository,
\texttt{docs/surcomplex/rank-one-berkovich/}, inspected 21 September 2026.
Already contains rank-one zero-free rigidity and a canonical product theorem;
its \texttt{prop:rankobstruction} is the rank-one boundary discussed in
\cref{ent:sec:collection}.

\bibitem{RepoDivisors}
\emph{Global Support Obstructions: Moving Divisors, Mittag--Leffler, and a
Picard Dichotomy}, research report in the \emph{Surreal} repository,
\texttt{docs/surcomplex/global-divisors/}, inspected 21 September 2026.
Stated over $\mathcal O(U)((t^\Gamma))$, a different coefficient ring.

\bibitem{RepoAnalyticGeometry}
\emph{Analytic Geometry over Surcomplex Hahn Fields}, research report in the
\emph{Surreal} repository, \texttt{docs/surcomplex/analytic-geometry/}.
Source of the ring-discipline convention invoked in \cref{ent:conv:ring}.

\bibitem{RepoAnalysis}
\emph{Surcomplex Analysis}, research report in the \emph{Surreal}
repository, \texttt{docs/surcomplex/analysis/}. Contains the all-scale
polynomial rigidity theorem discussed in \cref{ent:rem:rigidity-credit}.

\bibitem{RepoFoundations}
\emph{Foundations}, research report in the \emph{Surreal} repository,
\texttt{docs/foundations-and-computation/foundations/}. Source of the
localization principle and of the sign-flipped pair
$\sum t^{\pm n^2}z^n$ discussed in
\cref{ent:rem:criterion-strength,ent:sec:collection}.
\end{thebibliography}
\endgroup
\end{document}
